\documentclass[9pt]{article}
\usepackage[a4paper,margin=0.25in]{geometry}
\usepackage{amsmath,amssymb,amsfonts,mathtools}
\setlength{\jot}{4pt} 
\usepackage{braket}
\usepackage{slashed}
\usepackage{cancel}


\usepackage{graphicx}
\usepackage{tikz}
\usepackage{pgfplots}


\usepackage{xcolor}
\usepackage{hyperref}
\usepackage{comment}


\newcommand{\CGCket}{|{\rm CGC}\rangle}
\newcommand{\CGCbra}{\langle {\rm CGC}|}

\def\wt#1{\widetilde{#1}}
\def\mf#1{\mathbb{#1}}

\newcommand{\A}[3]{\mf{A}^{#2}_{#3}(#1)}

\numberwithin{equation}{section}

\title{Finite terms of the final cross section ($\Lambda\to0$, $\vee\to\infty$)}
\author{}
\date{}

\begin{document}
\maketitle
\noindent\textbf{What this file contains.} The NLO cross section is
$\mathrm{LO}+\log\frac{\vee}{k^{+}}\times[\text{BFKL formula}]+\log\frac{k^{+}}{\Lambda}\times[\text{JIMWLK formula}]+\text{DGLAP}+\text{RC}+\text{finite terms}$.
This file gives the finite terms row by row, and the DGLAP and running-coupling terms in their convergent form (Group II, Row IV), with $\Lambda\to0$ and $\vee\to\infty$ taken. The BFKL and JIMWLK formulas are those of the checked evolution files and are not repeated. Four-$\rho$ rows contribute only to the evolution terms.

\noindent\textbf{Conventions.} Phase $e^{-ik\cdot(w-w')}$; c.c.\ implied for Groups I and III; $\rho$ are non-commuting operators in the order written; principal value, logs are $\log|\dots|$; no $1/\epsilon$ poles are written (rule $\int d^{2}p/p^{2}=0$); one scale $\mu$. The soft plus-prescription terms are written with the recoil phase kept in the subtraction, so the JIMWLK terms stay exactly as in the JIMWLK formula and the subtraction generates the explicit extra finite terms shown. The $\log(1-k^{+}/\vee)$ terms go with the JIMWLK structure of their row and are not listed.

\noindent\textbf{Notation.} The two parts of the file were derived separately and use different auxiliary functions: Group I uses $\mathcal{E}(s)$ and $\mathcal{J}(s)$ defined in its first subsection ($p^{+}$ up to $k^{+}$); Groups II and III use $\widehat{\mathcal{E}}(s)$, $\widehat{\mathcal{J}}(s)$ and $\Psi(s)$ defined there ($p^{+}$ up to infinity). At large $|s|$ both $\mathcal{E}(s)$ and $\Psi(s)$ tend to $-\gamma_{E}-\log|s|-\frac{i\pi}{2}\,{\rm sgn}\,s$; this common tail function is what enters the cancellation of the large-distance tails between the rows.

\noindent{\color{magenta}\textbf{What is not finite or not settled} is marked \fbox{OPEN} in magenta (OPEN I-1 to I-3 in Group I, OPEN II/III-1 to II/III-4 in Groups II and III). In short: (1) the recoil terms ($\mathcal{E}$, $\Psi$) are not integrable over the soft position one by one; their $1/z^{2}$ tails cancel only in the sum over Group I Rows III, V and Group III Rows III, V, VI, under one integral; (2) the parts of these tails linear in $U(z)$ cancel only between the $w$ side and the $w'$ side, which is uniform only if $U(z)$ has a direction-independent limit at large distance; (3) Group II Row IV keeps a large-distance tail $\propto[-\frac{11}{6}+\frac16\cos2\varphi]$; (4) the combined kernel of Group I Row III converges at large $z$ only conditionally (angular integral first). No complete numerical integration over the transverse plane with Wilson lines has been done; the checks are pointwise limits, integrals with test functions, scaling exponents and tails.}

\noindent{\color{gray}\small Grey lines under the equations say from which checked equation each term comes and what was verified numerically.}

\section{Group I (with c.c.)}
\noindent\textbf{What is listed.} For every row the limit $\Lambda\to0$, $\vee\to\infty$ of
\begin{align}
\big[\text{checked row with }\Lambda,\vee\big]-\log\frac{\vee}{k^{+}}\times\big[\text{BFKL-list term of the row}\big]-\log\frac{k^{+}}{\Lambda}\times\big[\text{JIMWLK-list term of the row}\big]\,.\notag
\end{align}
The terms $\log\big(1-\frac{k^{+}}{\vee}\big)\times$(JIMWLK structure of the row) go with the JIMWLK term and are not listed. No $1/\epsilon$ pole is written ($\int d^{2}p/p^{2}=0$). All logarithms are $\log|\dots|$. ``c.c.'' is implied for every equation of this section.

\noindent\textbf{Notation.} $K^{i}(v)\equiv\frac{v^{i}}{v^{2}}$, $\;K(a,b,z)\equiv K(a-z)\cdot K(b-z)$, $\;L\equiv\log\frac{k^{+}}{\Lambda}$, and
\begin{align}
\mathcal{L}^{a}_{x'w'} &\equiv \big[U^{ab'}(x')-U^{ab'}(w')\big]\,\rho^{b'}(x')\,,\notag\\
\mathcal{C}_{\rm III}(z) &\equiv \mathcal{L}^{a}_{x'w'}\times\left[U^{de}(x)\,U^{bc}(z)\,\rho^{e}(x)\,\rho^{c}(y)-U^{de}(x)\,U^{bc}(y)\,\rho^{e}(x)\,\rho^{c}(y)\right]\qquad(\text{contracted with } f^{abd})\,,\notag\\
\mathcal{C}_{\rm IV} &\equiv \mathcal{L}^{a}_{x'w'}\left[f^{cdc'}\,U^{ac}(w)\,U^{bd}(z)\,U^{be}(y)\,\rho^{e}(y)\rho^{c'}(x)-f^{abc}\,U^{be}(y)\,U^{cd}(x)\,\rho^{e}(y)\rho^{d}(x)\right]\,.\notag
\end{align}
Two functions of one real variable appear in the terms that come from the plus prescriptions ($\mathrm{Ci}$, $\mathrm{Si}$: cosine and sine integrals):
\begin{align}
\mathcal{E}(s)&\equiv\int_{0}^{1}\frac{dt}{t}\,\big(e^{-ist}-1\big)=\mathrm{Ci}(|s|)-\gamma_{E}-\log|s|-i\,\mathrm{Si}(s)\,,\qquad
\mathcal{J}(s)\equiv\int_{0}^{1}dt\,e^{-ist}=\frac{1-e^{-is}}{is}\,,\label{eq:EJ}\\
\mathcal{E}(s)&\xrightarrow[s\to0]{}-is\,,\qquad \mathcal{E}(s)\xrightarrow[|s|\to\infty]{}-\gamma_{E}-\log|s|-\frac{i\pi}{2}\,{\rm sgn}(s)+{\cal O}(1/|s|)\,.\notag
\end{align}
{\color{gray}\small [Checked numerically: $\int_{\lambda}^{1-\lambda}\frac{dt}{t}e^{-ist}-\log\frac1\lambda-\mathcal{E}(s)=1.1\times10^{-2},\,1.1\times10^{-4},\,1.1\times10^{-6}$ for $s=0.3$ and $4.1\times10^{-1},\,4.1\times10^{-3},\,4.1\times10^{-5}$ for $s=40$ ($\lambda=10^{-2},10^{-4},10^{-6}$): the limit is pointwise in $s$, the error is ${\cal O}(\lambda s)$; large $s$: the remainder is $1/|s|$.]}

\noindent\textbf{How the soft plus-prescription terms are written (form (2): recoil phase kept in the subtraction).} A term $\int_{\Lambda}\frac{dp^{+}}{p^{+}}f(p^{+})$ of a checked row is written as
\begin{align}
\int_{\Lambda}^{k^{+}}\frac{dp^{+}}{p^{+}}\,f(p^{+})=\underbrace{L\,f(0)}_{\text{JIMWLK list, unchanged}}+\underbrace{\int_{0}^{k^{+}}\frac{dp^{+}}{p^{+}}\Big[f(p^{+})-f(0)\,e^{-i\frac{p^{+}}{k^{+}}s}\Big]}_{\text{regular term}}+\underbrace{f(0)\,\mathcal{E}(s)}_{\text{recoil term}}\,,\label{eq:form2}
\end{align}
$s=k\cdot(\text{soft position}-\text{hard position})$, where ``hard position'' is the transverse position of the measured gluon in $f(0)$ (the point called $w$ in the JIMWLK formula) and ``soft position'' is the position of the gluon whose $p^{+}\to0$ (called $z$ there). The regular terms below are integrable at large soft position by themselves (checked). The recoil terms are not: each behaves like $\frac{T}{z^{2}}\,\mathcal{E}(k\cdot z)$ at large soft position $z$, with $T$ the $1/z^{2}$ tail of $f(0)$; only their sum over rows is integrable (Sec.~\ref{sec:tailsI}, {\color{magenta}OPEN I-2}).

\subsection{Row I}
\noindent No term (Row I contributes only to the four-$\rho$ terms).

\subsection{Row II}
\begin{align}
\frac{d^{3}N_{1}^{\rm II}}{d^{3}k}\bigg|_{\rm finite} & = \frac{1}{(2\pi)^{3}}\, \frac{g^{4}\,N_{c}}{32\pi^{4}}\,\frac{1}{k^{+}}\,\left\{-\frac{67}{9}+\frac{\pi^{2}}{3}\right\} \int_{w, w',x,x'} e^{-ik\cdot(w-w^\prime)}\, K(x'-w')\cdot K(x-w)\;\mathcal{L}^{a}_{x'w'}\,\left[U^{ab}(w) - U^{ab}(x)\right]\rho^{b}(x)\,.\label{eq:IIa}
\end{align}
{\color{gray}\small [From $N_{1}$ of the checked Row II and $N_{2b\,\rm UV}$ of the checked Row V. With $\mathcal{B}\equiv2\gamma_{E}+\log\frac{(x-w)^{2}\mu^{2}}{4}$ and the relabelling $x\to w$, $y\to x$ in Row V, the two rows are $\big\{-\mathcal{B}\big[\frac{11}{3}-2L+\log(1-\frac{k^{+}}{\vee})\big]-\frac{67}{9}+\frac{\pi^{2}}{3}+\log^{2}(1-\frac{k^{+}}{\vee})+\log(1-\frac{k^{+}}{\vee})\log\frac{\vee}{k^{+}}\big\}$ and $\big\{+\mathcal{B}\big[\frac{11}{3}-4L\big]\big\}$ times the same integrand. The bracket $\big[2\gamma_{E}+\log\frac{(x-w)^{2}\mu^{2}}{4}\big]\times\frac{11}{3}$ of Row II and the bracket $\big[2\gamma_{E}+\log\frac{(y-x)^{2}\bar\mu^{2}}{4}\big]\times\big(-\frac{11}{6}\big)$ of Row V ($N_{2b\,\rm UV}$) cancel exactly ($\bar\mu=\mu$; $-\frac{11}{3}+2\cdot\frac{11}{6}=0$): \emph{neither appears in the final result, and no running-coupling term is left in Group I.} The terms $+2L\,\mathcal{B}$, $-4L\,\mathcal{B}$ are the JIMWLK-list terms of the two rows. Numerically, row II $+$ row V $-$ JIMWLK terms $-\big(-\frac{67}{9}+\frac{\pi^{2}}{3}\big)=-7\times10^{-2},\,-9\times10^{-4},\,-1.4\times10^{-5}$ for $\frac{k^{+}}{\vee}=10^{-2},10^{-4},10^{-6}$ (it is $-\mathcal{B}\log(1-\frac{k^{+}}{\vee})+\log^{2}(1-\frac{k^{+}}{\vee})+\log(1-\frac{k^{+}}{\vee})\log\frac{\vee}{k^{+}}$, a number times the LO-type integrand).]}
\begin{align}
\frac{d^{3}N_{2b}^{\rm II}}{d^{3}k} & = -\frac{1}{(2\pi)^{3}}\,\frac{g^{4}}{64\pi^{6}}\,\frac{1}{k^{+}}\int_{w,w',x'} e^{-ik\cdot(w-w')}
\int d^{2}p\,d^{2}q\,\int_{x,y} e^{iq\cdot(w-y)}\,e^{ip\cdot(y-x)}\,K^{i}(x'-w')\,\frac{1}{\mathbf{q}^{2}}\,\mathcal{R}^{i}(\mathbf{q},\mathbf{p})\notag\\
&\qquad\times\mathcal{L}^{a}_{x'w'}\,
\Big[f^{hbc}\,U^{ah}(w)\left\{\rho^{b}(x),\rho^{c}(y)\right\}-f^{abc}\,U^{bd}(x)\,U^{ce}(y)\left\{\rho^{d}(x),\rho^{e}(y)\right\}\Big]\,,\label{eq:IIb}\\
\mathcal{R}^{i}(\mathbf{q},\mathbf{p})&\equiv\frac12\left[\frac{R^{i}_{0}(\mathbf{q},\mathbf{p})}{\mathbf{p}^{2}}-\frac{R^{i}_{0}(\mathbf{q},\mathbf{q}-\mathbf{p})}{(\mathbf{q}-\mathbf{p})^{2}}\right],\qquad
R^{i}_{0}(\mathbf{q},\mathbf{p}) = -2q^{i}\,\log\frac{p^{2}}{q^{2}}
+\frac{(2p^{2}-2\,q\cdot p+q^{2})\,q^{i}+q^{2}p^{i}}{(q-p)^{2}}\,
\log\frac{q^{2}}{(q-p)^{2}}\,.\notag
\end{align}
{\color{gray}\small [From $N_{2b}$ of the checked Row II. \emph{Same kernel for both colour structures:} for the barred structure the row has $R^{i}$ (no $\log\frac{\vee}{k^{+}}$); for the unbarred structure (operator ordering $A^{(1)}N_{2}$) the row has $R^{i}+2q^{i}\log\frac{\vee}{k^{+}}$, and $2q^{i}\log\frac{\vee}{k^{+}}$ times the unbarred structure is exactly the Row II term of the BFKL list. After it is taken out, the unbarred structure has $R^{i}$ as well and no finite remainder is left from the ordering term: at a random point, $[R^{i}+2q^{i}\log\frac{\vee}{k^{+}}]-\log\frac{\vee}{k^{+}}\,2q^{i}-R^{i}_{0}=(-0.27,-0.12)$, $(-2.8,-1.3)\times10^{-3}$, $(-2.8,-1.3)\times10^{-5}$, $(-2.8,-1.3)\times10^{-7}$ for $\frac{k^{+}}{\vee}=10^{-2},10^{-4},10^{-6},10^{-8}$, identical to $R^{i}-R^{i}_{0}$ of the barred structure. \emph{Infrared points:} $R^{i}_{0}/(p^{2}q^{2})$ alone is not integrable at $p\to0$ and at $q\to p$ ($\log p^{2}/p^{2}$). Both colour structures are antisymmetric under $x\leftrightarrow y$, i.e. under $\mathbf{p}\to\mathbf{q}-\mathbf{p}$; in the antisymmetrised kernel $\mathcal{R}^{i}$ the two points cancel each other, and this is the value that the rule $\int d^{2}p/p^{2}=0$ assigns (the rule allows the relabelling). Numerically $|\mathcal{R}|/q^{2}\simeq3.5\,\log(1/|p|)/|p|$ at $p\to0$ (ratio $2.97,\,3.30,\,3.47,\,3.57$ for $|p|=10^{-2}\dots10^{-5}$), $\int_{|p|<\delta}d^{2}p\,|\mathcal{R}|/q^{2}=5.6,\,1.3,\,0.33$ for $\delta=10^{-1},10^{-2},10^{-3}$; $q\to0$: integrable; $|p|\to\infty$: $|p|^{-3}$; $|q|\to\infty$: $\log(q^{2})/q^{2}$, convergent through the phase $e^{iq\cdot(w-y)}$ as at LO. $\vee\to\infty$ integrated over $p$ with $e^{-p^{2}}$: difference $-0.20,\,-2.4\times10^{-2},\,-2.5\times10^{-3},\,-2.5\times10^{-4}$ for $\frac{k^{+}}{\vee}=10^{-1}\dots10^{-4}$.]}

\subsection{Row III}
\noindent Common prefactor and colour factor. $\mathcal{C}_{\rm III}$ does not contain $U(w)$.
\begin{align}
\frac{d^{3}N^{\rm III}}{d^{3}k}\bigg|_{\rm finite} & =\frac{1}{(2\pi)^{3}}\,\frac{ig^{4}f^{abd}}{8\pi^{5}}\,\frac{1}{k^{+}}\int_{w,w'}e^{-ik\cdot(w-w')}\int_{x,y,z,x'}\Bigg\{
K^{i}(x'-w')\,K^{j}(y-z)\;\mathcal{K}^{ij}\;\mathcal{C}_{\rm III}(z)\notag\\
&\quad-\frac12\,\log\left[\frac{(z-w)^{2}}{(x-w)^{2}}\right]K(y,w,z)\;K(x'-w')\cdot K(x-w)\;\mathcal{C}_{\rm III}(z)\notag\\
&\quad-\,\mathcal{E}\big(k\cdot(z-w)\big)\,K(x,w,z)\;K(x'-w')\cdot K(y-w)\;\mathcal{C}_{\rm III}(w)\Bigg\}\,,\label{eq:III}
\end{align}
where $\mathcal{C}_{\rm III}(w)$ is $\mathcal{C}_{\rm III}$ with $U^{bc}(z)\to U^{bc}(w)$, and
\begin{align}
\mathcal{K}^{ij}&=\log\left[\frac{(x-z)^{2}}{(x-w)^{2}}\right]\Bigg\{\frac{\delta^{ij}}{2\,(z-w)^{2}}-\frac{(x-w)^{i}\,(x-z)^{j}}{(x-w)^{2}\,(x-z)^{2}}
+\frac{1}{(z-w)^{2}\left[(x-w)^{2}-(x-z)^{2}\right]}\Big[(x-w)^{i}(z-w)^{j}\,\frac{(x-z)^{2}}{(x-w)^{2}}\notag\\
&\qquad\quad+(z-w)^{i}(x-z)^{j}\,\frac{(x-w)^{2}}{(x-z)^{2}}-(z-w)^{i}(x-w)^{j}-(x-z)^{i}(z-w)^{j}\Big]\Bigg\}\notag\\
&\quad+\frac{\mathcal{A}_{0}}{(z-w)^{2}}\Bigg\{(z-w)^{i}(x-w)^{j}+(x-z)^{i}(z-w)^{j}-(x-w)\cdot(x-z)\left[\frac{(z-w)^{i}(x-z)^{j}}{(x-z)^{2}}+\frac{(x-w)^{i}(z-w)^{j}}{(x-w)^{2}}\right]\Bigg\}\notag\\
&\quad-\frac{1}{2\,(z-w)^{2}}\log\left[\frac{(z-w)^{2}}{(x-w)^{2}}\right]\frac{(z-w)^{i}(x-z)^{j}}{(x-z)^{2}}\,,\label{eq:Kij}\\
\mathcal{A}_{0} & = \frac{1}{\sqrt{(x-w)^{2}\,(x-z)^{2}-\left[(x-w)\cdot (x-z)\right]^{2}}}\,\mathrm{arc\,cos}\,\left(\frac{(x-w)\cdot (z-w)}{\sqrt{(x-w)^{2}\,(z-w)^{2}}}\right)\,.\notag
\end{align}
{\color{gray}\small [From $N_{1}$, $N_{2}$, $N_{3a}$, $N_{3b}$ of the checked Row III with $\Lambda$, $\vee$ kept (the $1/\Lambda$, $1/(\vee-k^{+})$ terms cancel). The pointwise coefficient of $L$ in the row is $\frac12K(x'-w')\cdot K(x-w)\,K(x,y,z)\,\mathcal{C}_{\rm III}(z)-\frac{K(x'-w')\cdot K(y-z)}{(z-w)^{2}}\,\mathcal{C}_{\rm III}(z)$ (checked numerically to $10^{-16}$); the first is the JIMWLK-list term, the second is treated in Eqs.~\eqref{eq:IIIA1}--\eqref{eq:IIIA3}. Pointwise check at random points: [row $-\log\frac{\vee}{k^{+}}\times$coeff.\ $-L\times$coeff.] $-$ [first two lines of \eqref{eq:III} $-\frac12\log\frac{(z-w)^{2}}{(x-w)^{2}}\frac{K(x'-w')\cdot K(y-z)}{(z-w)^{2}}$] $=1.0\times10^{-3},\,3.9\times10^{-4},\,4.1\times10^{-5},\,4.2\times10^{-6},\,4.2\times10^{-7}$ for $\frac{\Lambda}{k^{+}}=\frac{k^{+}}{\vee}=10^{-2}\dots10^{-6}$ (second point: $8.3\times10^{-4},\,8.3\times10^{-6},\,8.3\times10^{-8}$ for $10^{-3},10^{-5},10^{-7}$).]}

\noindent\textbf{First line of \eqref{eq:III}} (the 24 terms that are neither of the two flagged terms): only the combined kernel is integrable. {\color{gray}\small [Checked for $\mathcal{K}^{ij}$: $z\to w$: $\simeq3\log(1/|z-w|)/|z-w|$, angular mean of the $1/(z-w)^{2}$ coefficient $\to0$ ($-1.4\times10^{-3}\dots-1.5\times10^{-8}$ for $|z-w|=10^{-2}\dots10^{-5}$); no pole on $(x-z)^{2}=(x-w)^{2}$ (values $0.19948358$ and $0.19948353$ at distance $\pm10^{-7}$ from the line); finite on the collinear configurations; $x\to w$: at most $1/|x-w|$; $x\to z$: regular. Integrated over $z$ for the $z$-independent part of $\mathcal{C}_{\rm III}$ (with $K^{j}(y-z)$, angular integral first): [24 terms of the row with $\Lambda$, $\vee$ minus their logarithms] $-$ [first line of \eqref{eq:III}] $=-1.66,\,-0.21,\,-2.8\times10^{-2},\,-3.0\times10^{-3}$ for $\frac{\Lambda}{k^{+}}=\frac{k^{+}}{\vee}=10^{-1}\dots10^{-4}$ (the last one known to $\pm2\times10^{-4}$).]} At large $z$ the kernel does \emph{not} decay like $|z|^{-3}$ pointwise (as stated in the limits file) but like
\begin{align}
K^{j}(y-z)\,\mathcal{K}^{ij}\;\xrightarrow[|z|\to\infty]{}\;\frac{\varphi}{z^{2}}\;\frac{e_{\perp}^{i}}{|x-w|}+{\cal O}\Big(\frac{\log|z|}{|z|^{3}}\Big)\,,\qquad -\pi<\varphi<\pi\,,\label{eq:saw}
\end{align}
$\varphi$ = angle from $x-w$ to $z$ (counterclockwise), $e_{\perp}$ = unit vector $\frac{x-w}{|x-w|}$ rotated by $+\frac{\pi}{2}$; it comes from the $\mathcal{A}_{0}$ terms. The coefficient is odd in $\varphi$: its angular average vanishes. {\color{gray}\small [Checked: $z^{2}\,K^{i}(x'-w')K^{j}(y-z)\mathcal{K}^{ij}=-0.264573,\,1.495434,\,-3.511430$ against $\varphi\,K(x'-w')\cdot e_{\perp}/|x-w|=-0.264373,\,1.495565,\,-3.511692$ at $|z|=10^{4}$, three directions; angular mean $\times z^{2}=-6.8\times10^{-2},\,-6.7\times10^{-3},\,-3.9\times10^{-4}$ at $|z|=10^{2},10^{3},10^{4}$. At finite $\Lambda$ the tail is $\pm\psi_{\rho}$ with $\psi_{\rho}=\arg\big(\rho+e^{i|\varphi|}\big)$, $\rho=\frac{\Lambda|z|}{k^{+}|x-w|}$: also odd in $\varphi$ and cut off at $|z|\sim\frac{k^{+}}{\Lambda}|x-w|$.]}

\noindent{\color{magenta}\fbox{OPEN I-1} The first line of \eqref{eq:III} is integrable at large $z$ only conditionally: for the part of $\mathcal{C}_{\rm III}$ without $U(z)$ the angular integral has to be done first (then the $z$ integral converges, and this is also what the row gives at finite $\Lambda$, where the tail is odd in $\varphi$ at every $|z|$); for the part with $U^{bc}(z)$ it converges only if $U(z)$ tends to a direction-independent limit at large $|z|$. Otherwise a term $\propto L\int d\varphi\,\varphi\,U(z\to\infty,\varphi)$ would be left. Nothing is assumed about $U$ at large distance, so this is not decided here.}

\noindent\textbf{Second line of \eqref{eq:III}} (third term of $N_{3b}$, FLAG I-4; $z$ is the position of the soft gluon). It is [LO kernel] $-$ [LO kernel with the recoil shift]: with $a_{\xi}\equiv(x-w)-\xi\,(x-z)$,
\begin{align}
-\frac12\log\left[\frac{(z-w)^{2}}{(x-w)^{2}}\right]K^{i}(x-w)=-\int_{0}^{1}\frac{d\xi}{\xi}\left[K^{i}(x-w)-K^{i}(x-w)\,\frac{(x-w)\cdot a_{\xi}}{a_{\xi}^{2}}\right]\,;\qquad\text{in the row:}\ \int_{\Lambda/k^{+}}^{1-\Lambda/k^{+}}\frac{d\xi}{\xi}\,.\label{eq:IIIB}
\end{align}
The subtraction is already of the recoil-shifted kind and nothing of this term is in the JIMWLK list (its pointwise coefficient of $L$ vanishes). It is integrable at $z\to w$ ($\log/|z-w|$), $z\to y$, $x\to w$. At large $z$ it behaves like $\frac{1}{z^{2}}\big[-\log\frac{|z|}{|x-w|}\big]K(x'-w')\cdot K(x-w)\,\mathcal{C}_{\rm III}(z)$, i.e. it is of the same kind as the recoil terms, with $-\log\frac{|z-w|}{|x-w|}$ in the place of $\mathcal{E}$ (Sec.~\ref{sec:tailsI}). {\color{gray}\small [Checked: integrated over $z$ with $(1+z^{2})^{-1/2}$ in the place of the colour structure, (row term at $\Lambda$) $-$ (second line) $=-0.24,\,-3.0\times10^{-2},\,-3.7\times10^{-3},\,-4.4\times10^{-4}$ for $\frac{\Lambda}{k^{+}}=10^{-1}\dots10^{-4}$ (value of the integral $1.034$); tail: $z^{2}\times$(angular mean)$/\log\frac{|z|}{|x-w|}=2.32427,\,2.32426,\,2.32426$ at $|z|=10^{2},10^{3},10^{4}$ against $-K(x'-w')\cdot K(x-w)=2.32426$.]}

\noindent\textbf{Third line of \eqref{eq:III}} (last term of $N_{3b}$ together with the term $-\big[L+\log\big(1-\frac{\Lambda}{k^{+}}\big)\big]\frac{K(x'-w')\cdot K(y-z)}{(z-w)^{2}}\mathcal{C}_{\rm III}(z)$ of $N_{3a}$; FLAG I-3). Derivation. With $\Phi(z)\equiv K(x'-w')\cdot K(y-z)\,\mathcal{C}_{\rm III}(z)$, $\lambda\equiv\frac{\Lambda}{k^{+}}$, the two terms of the row are (in units of the prefactor of \eqref{eq:III}, exactly)
\begin{align}
-\frac12\int_{z}\frac{\mathcal{G}_{\Lambda}}{(z-w)^{2}}\,\Phi(z)\,,\qquad
\mathcal{G}_{\Lambda}&=\log\frac{\big[(x-w)+\frac{1-\lambda}{\lambda}(z-w)\big]^{2}}{\big[(x-w)+\frac{\lambda}{1-\lambda}(z-w)\big]^{2}}
=\int_{\frac{\lambda}{1-\lambda}}^{\frac{1-\lambda}{\lambda}}\frac{d\xi}{\xi^{2}}\,\frac{2\,\big[(x-w)+\frac{z-w}{\xi}\big]\cdot(z-w)}{\big[(x-w)+\frac{z-w}{\xi}\big]^{2}}\,.\label{eq:IIIA1}
\end{align}
Substituting $z=w-\xi(z_{s}-w)$ ($z_{s}$: position of the soft gluon, $z$: position of the hard gluon before it absorbs the soft one),
\begin{align}
-\frac12\int_{z}\frac{\mathcal{G}_{\Lambda}}{(z-w)^{2}}\,\Phi(z)=-\int_{\frac{\lambda}{1-\lambda}}^{\frac{1-\lambda}{\lambda}}\frac{d\xi}{\xi}\int_{z_{s}}K(x,w,z_{s})\;\Phi\big(w-\xi(z_{s}-w)\big)\,,\label{eq:IIIA2}
\end{align}
and, under $\int_{w}e^{-ik\cdot w}$, changing the variable $w$ to the hard position $w_{h}=w-\xi(z_{s}-w)$ (then $K(x,w,z_{s})=(1+\xi)K(x,w_{h},z_{s})$, $d^{2}w=\frac{d^{2}w_{h}}{(1+\xi)^{2}}$, $e^{-ik\cdot w}=e^{-ik\cdot w_{h}}e^{-i\frac{\xi}{1+\xi}k\cdot(z_{s}-w_{h})}$) and $t=\frac{\xi}{1+\xi}$,
\begin{align}
\int_{w}e^{-ik\cdot w}\Big(-\frac12\Big)\int_{z}\frac{\mathcal{G}_{\Lambda}}{(z-w)^{2}}\,\Phi(z)
&=-\int_{w,z}e^{-ik\cdot w}\,K(x,w,z)\,\Phi(w)\int_{\lambda}^{1-\lambda}\frac{dt}{t}\,e^{-it\,k\cdot(z-w)}\notag\\
&=-\int_{w,z}e^{-ik\cdot w}\,K(x,w,z)\,\Phi(w)\,\Big[L+\mathcal{E}\big(k\cdot(z-w)\big)+{\cal O}(\lambda)\Big]\,.\label{eq:IIIA3}
\end{align}
This identity is exact in $\lambda$ (first line). The term with $L$ is the JIMWLK-list term in the author's form (kernel $-K(x'-w')\cdot K(y-w)\,K(x,w,z)$, Wilson line $U(w)$); the rest is the third line of \eqref{eq:III}. It is a pure recoil term: $f(p^{+})=f(0)\,e^{-i\frac{p^{+}}{k^{+}}k\cdot(z-w)}$, no regular term. Integrable at $z\to w$ ($\mathcal{E}\to0$), $z\to x$, $y\to w$; large $z$: $\frac{1}{z^{2}}\mathcal{E}(k\cdot z)$ (Sec.~\ref{sec:tailsI}). {\color{gray}\small [Checked: \eqref{eq:IIIA1} at a random point for $\lambda=0.3,\,0.1,\,10^{-2},\,10^{-3}$ (10 digits). The first line of \eqref{eq:IIIA3} with a Gaussian regulator, i.e. $-\frac12\int d^{2}\zeta\,e^{ik\cdot\zeta}e^{-\delta\zeta^{2}}\zeta^{-2}\big[\log(V+\frac{\zeta}{\lambda})^{2}-\log(V+\frac{\zeta}{1-\lambda})^{2}\big]=-\int_{\lambda}^{1-\lambda}\frac{dt}{t}\int d^{2}\zeta'\,e^{-itk\cdot\zeta'}e^{-\delta t^{2}\zeta'^{2}}\frac{(\zeta'-V)\cdot\zeta'}{(\zeta'-V)^{2}\zeta'^{2}}$ ($V=x-z$ resp. $x-w$; $\delta=0.05$): $-10.896402+0.822556\,i$ on both sides for $\lambda=0.2$ (difference $3\times10^{-13}$), $-34.188243+1.392118\,i$ for $\lambda=0.05$ (difference $3\times10^{-10}$). The range of $\xi$ is $\frac{\lambda}{1-\lambda}<\xi<\frac{1-\lambda}{\lambda}$, i.e. $\Lambda<q^{+}<k^{+}-\Lambda$ for $q^{+}=t\,k^{+}$, not $\Lambda<q^{+}<k^{+}$ in the variable $\xi$ as written in the note on the flagged terms. Without a regulator, [row] $-L\times$[JIMWLK-list term] has no limit for this term alone after the $z$ integration, because the JIMWLK-list term itself has the tail $-\frac{1}{z^{2}}\Phi(w)$: the statement is the pointwise one, $\int_{\lambda}^{1-\lambda}\frac{dt}{t}e^{-its}-L\to\mathcal{E}(s)$, and the tail is dealt with in Sec.~\ref{sec:tailsI}.]}

\subsection{Row IV}
\begin{align}
\frac{d^{3}N^{\rm IV}}{d^{3}k}\bigg|_{\rm finite} & = \frac{1}{(2\pi)^{3}}\,\frac{ig^{4}}{8\pi^{5}}\,\frac{1}{k^{+}}
\int_{x,y,z,w,x',w'}e^{-ik\cdot (w-w^\prime)}\,K^{i}(x'-w')\,K^{j}(y-z)\,\log\left[\frac{(x-z)^{2}}{(x-w)^{2}}\right]\notag\\
&\quad\times\bigg[K^{j}(x-z)\,K^{i}(z-w)-K^{i}(x-w)\,K^{j}(z-w)-K^{j}(x-z)\,K^{i}(x-w)+\frac{\delta^{ij}}{2\,(z-w)^{2}}\bigg]\;\mathcal{C}_{\rm IV}\,.\label{eq:IV}
\end{align}
{\color{gray}\small [From the checked Row IV: the row is $\big[\log\frac{\vee}{k^{+}}+\log\frac{(1-\frac{k^{+}}{\vee})(x-z)^{2}+\frac{k^{+}}{\vee}(x-w)^{2}}{(x-w)^{2}+\frac{\Lambda}{k^{+}}(x-z)^{2}}\big]\times$(bracket) $-\frac{K(x'-w')\cdot K(y-z)}{2(z-w)^{2}}\big[\log\frac{\vee}{k^{+}}-\log(1+\frac{\Lambda}{k^{+}})\big]+\big[L+\log\frac{\vee}{k^{+}}+\log(1-\frac{k^{+}}{\vee})\big]\times$(JIMWLK structure). The term $\log\big(1-\frac{k^{+}}{\vee}\big)\times K^{i}(x'-w')K^{j}(y-z)\big[K^{i}(x-w)K^{j}(z-w)+\frac12K^{j}(x-z)K^{i}(x-w)\big]\mathcal{C}_{\rm IV}$ is dropped: it goes with the JIMWLK term of the row. Pointwise: row $-$ logarithms $-$ \eqref{eq:IV} $=$ $2.2\times10^{-3},\,2.2\times10^{-5},\,2.2\times10^{-7}$ for $\frac{\Lambda}{k^{+}}=\frac{k^{+}}{\vee}=10^{-2},10^{-4},10^{-6}$. Integrated over $z$ for the $z$-independent part of $\mathcal{C}_{\rm IV}$ (with $K^{j}(y-z)$, angular integral first): $-0.49,\,-9.2\times10^{-2},\,-1.4\times10^{-2},\,-1.8\times10^{-3}$ for $10^{-1}\dots10^{-4}$. Integrability: $z\to w$: $1/|z-w|$ (the logarithm vanishes linearly; the $\delta^{ij}$ term and the $-\frac12[\dots]$ term of the row must stay together); $x\to z$, $x\to w$: $\log/|\cdot|$; large $z$: $z^{2}\times$kernel$\times K^{j}(y-z)\to0$ like $\log|z|/|z|$ pointwise (maximum over the direction $3.1\times10^{-2},\,4\times10^{-3},\,6\times10^{-4},\,1\times10^{-4}$ at $|z|=10^{2}\dots10^{5}$): no $1/z^{2}$ tail, in contrast to Row III.]}

\noindent{\color{magenta}\fbox{OPEN I-3} The dropped $\log\big(1-\frac{k^{+}}{\vee}\big)$ term of Row IV multiplies a $z$ integral that diverges logarithmically at large $z$ for the $z$-independent part of $\mathcal{C}_{\rm IV}$ (tail $-\frac{1}{2z^{2}}K(x'-w')\cdot K(x-w)$). It is the same integral as in the $\log\frac{k^{+}}{\Lambda}$ and $\log\frac{\vee}{k^{+}}$ terms of the row; whether it is finite is a question about the JIMWLK and BFKL formulae (non-singlet $\int d^{2}z/z^{2}$ term), not about the finite terms. It vanishes for $\vee\to\infty$ only if that integral is regulated first.}

\subsection{Row V}
\noindent Variables of this row: $x$ is the position of the gluon with $p^{+}=\xi k^{+}$, $z$ the position of the gluon with $k^{+}-p^{+}$; the $w$ integral is done, $w=z+\xi(x-z)$. Throughout
\begin{align}
s\equiv k\cdot(x-z)\,,\qquad e^{-ik\cdot(x-w')}=e^{-ik\cdot(z-w')}\,e^{-is}\,.\notag
\end{align}
For $p^{+}\to0$ the soft gluon is at $x$ and the hard one at $z$ (recoil phase $e^{-i\xi s}$); for $p^{+}\to k^{+}$ the soft gluon is at $z$ and the hard one at $x$ (recoil phase $e^{-i(1-\xi)k\cdot(z-x)}$).
\subsubsection*{First part}
\begin{align}
\frac{d^{3}N_{1}^{\rm V}}{d^{3}k}\bigg|_{\rm finite}  & =  \frac{1}{(2\pi)^{3}}\, \frac{ig^{4}}{32\pi^{5}}\, \frac{1}{k^{+}}\,\int_{x,z,y,y',x',w'}\,e^{-ik\cdot (z-w^\prime)}\,\Big\{\mathcal{J}(s)\;K(x'-w')\cdot K(z-x)\;K(y-z)\cdot K(y'-x)\notag \\
&\qquad-\,\mathcal{E}(s)\;K(x'-w')\cdot K(y-z)\;K(z-x)\cdot K(y'-x)
-\,e^{-is}\,\mathcal{E}(-s)\;K(x'-w')\cdot K(y'-x)\;K(z-x)\cdot K(y-z)\Big\}\notag \\
&\qquad\times\,\mathcal{L}^{a}_{x'w'}\, \bigg[f^{adc}\,U^{db}(x)\,U^{ce}(z)\, \{\rho^{e}(y),\rho^{b}(y')\}-f^{abc}\,U^{cd}(y)\,U^{be}(y')\, \{\rho^{d}(y),\rho^{e}(y')\}\bigg]\,.\label{eq:V1}
\end{align}
{\color{gray}\small [From the first part of the checked Row V with $\Lambda<p^{+}<k^{+}-\Lambda$: bracket $\delta_{im}\delta_{jk}-\frac{k^{+}}{p^{+}}\delta_{jm}\delta_{ik}-\frac{k^{+}}{k^{+}-p^{+}}\delta_{ij}\delta_{km}$ times the phase $e^{-i\xi s}$; the transverse kernels do not depend on $p^{+}$, so $f(p^{+})=f(0)e^{-i\xi s}$ exactly and \eqref{eq:form2} has no regular term from the two plus prescriptions: they are replaced by the two recoil terms (second and third term). The first term is the $\delta_{im}\delta_{jk}$ term. Numerically $\int_{\lambda}^{1-\lambda}d\xi\,[\text{row}]-L\,[\text{JIMWLK coefficients}]-\{\dots\}=1.6\times10^{-2},\,1.6\times10^{-4},\,1.6\times10^{-6}$ for $\lambda=10^{-2},10^{-4},10^{-6}$ at a random point. Large distance of the $\mathcal{J}$ term: $R^{2}\times$(angular mean) $=1.4\times10^{-4},\,3.4\times10^{-5},\,3.8\times10^{-7}$ ($x$ far), $8.9\times10^{-3},\,5.0\times10^{-4},\,2.9\times10^{-5}$ ($z$ far) at $R=30,\,300,\,3000$.]}
\subsubsection*{Second part}
\noindent $N_{2a}$ and $N_{2b\,\rm rest}$ have the same kernel. With
\begin{align}
B_{1}^{im}&\equiv\frac{(y-x)^{i}\,(x-z)^{m}}{(x-z)^{2}}-\frac{(y-x)^{i}\,(y-z)^{m}}{2\,(y-z)^{2}}\,,\qquad
B_{2}^{im}\equiv\frac{(y-z)^{i}\,(x-z)^{m}}{(x-z)^{2}}+\frac{(y-x)^{m}\,(y-z)^{i}}{2\,(y-x)^{2}}\,,\notag\\
T_{1}&\equiv\frac{(y-z)\cdot (x-z)}{(x-z)^{2}}+\frac{(y-x)\cdot (y-z)}{2\,(y-x)^{2}}-\frac{(y-x)^{2}-(y-z)^{2}}{2\,(x-z)^{2}}\,,\notag\\
T_{2}&\equiv\frac{(y-x)\cdot (x-z)}{(x-z)^{2}}-\frac{(y-x)\cdot (y-z)}{2\,(y-z)^{2}}-\frac{(y-x)^{2}-(y-z)^{2}}{2\,(x-z)^{2}}\,,\notag\\
Q^{im}&\equiv\frac{(y-x)^{m}\,(x-z)^{i}}{(x-z)^{2}}-\frac{(y-x)^{m}\,(y-z)^{i}}{2\,(y-z)^{2}}+\frac{(y-z)^{m}\,(x-z)^{i}}{(x-z)^{2}}+\frac{(y-x)^{i}\,(y-z)^{m}}{2\,(y-x)^{2}}\,,\notag
\end{align}
the kernel is
\begin{align}
\mathcal{F}^{im}&\equiv\int_{0}^{1}d\xi\;\frac{e^{-i\xi s}}{\xi\,(y-x)^{2}+(1-\xi)\,(y-z)^{2}}\;\mathcal{N}^{im}(\xi)\;-\;\mathcal{E}(s)\,\frac{B_{2}^{im}}{(y-z)^{2}}\;-\;e^{-is}\,\mathcal{E}(-s)\,\frac{B_{1}^{im}}{(y-x)^{2}}\,,\label{eq:F}\\
\mathcal{N}^{im}(\xi)&\equiv\delta^{im}\,\xi(1-\xi)\,\frac{(y-x)^{2}-(y-z)^{2}}{(x-z)^{2}}+\delta^{im}\big[(1-\xi)\,T_{1}+\xi\,T_{2}\big]+\frac{(y-x)^{2}}{(y-z)^{2}}\,B_{2}^{im}+\frac{(y-z)^{2}}{(y-x)^{2}}\,B_{1}^{im}-Q^{im}\,,\notag
\end{align}
(first term of $\mathcal{F}$: regular term; second and third: recoil terms, soft gluon at $x$ resp. at $z$), and
\begin{align}
\frac{d^{3}N_{2a}^{\rm V}}{d^{3}k}\bigg|_{\rm finite}  & =  \frac{1}{(2\pi)^{3}}\, \frac{g^{4}}{16\pi^{5}}\,\frac{1}{k^{+}}\,\int_{x,y,z,x',w'}\, e^{-ik\cdot(z-w')}\,K^{i}(x'-w')\,K^{m}(z-x)\;\mathcal{F}^{im}\notag \\
&\qquad\times \Big[U^{ab'}(x')-U^{ab'}(w')\Big]\,\Big[-f^{adc}\,f^{fbe}\,U^{db}(x)\,U^{ce}(z)+N_{c}\,U^{af}(x)\Big]\,\rho^{b'}(x')\,\rho^{f}(y)\,,\label{eq:V2a}\\
\frac{d^{3}N_{2b\,\rm rest}^{\rm V}}{d^{3}k}\bigg|_{\rm finite}  & =  \frac{1}{(2\pi)^{3}}\, \frac{g^{4}\,N_{c}}{16\pi^{5}}\,\frac{1}{k^{+}}\,\int_{x,y,z,x',w'}\, e^{-ik\cdot(z-w')}\,\Bigg\{K^{i}(x'-w')\,K^{m}(z-x)\;\mathcal{F}^{im}\notag\\
&\qquad+\frac{11}{6}\;e^{-is}\;K(x'-w')\cdot K(y-x)\;\frac{K(y-z)\cdot(y-x)}{(x-z)^{2}}\Bigg\}\notag\\
&\qquad\times\Big[U^{ab'}(x')-U^{ab'}(w')\Big]\,\Big[U^{af}(y)-U^{af}(x)\Big]\rho^{b'}(x')\,\rho^{f}(y)\,.\label{eq:V2b}
\end{align}
$N_{2b\,\rm UV}$ of Row V is cancelled by Row II (note under Eq.~\eqref{eq:IIa}). The two colour structures add up to $[U(x')-U(w')]^{ab'}\big[-f^{adc}f^{fbe}U^{db}(x)U^{ce}(z)+N_{c}U^{af}(y)\big]\rho^{b'}(x')\rho^{f}(y)$, the single colour structure of the checked row before it was split.

\noindent{\color{gray}\small [From the second part of the checked Row V with $\Lambda<p^{+}<k^{+}-\Lambda$ (one $\xi$ integral, bracket with $-\frac{\xi}{1-\xi}B_{1}$ and $-\frac{1-\xi}{\xi}B_{2}$), $d_{\perp}\to2$. Derivation: $\frac{1}{\xi}\big[\frac{1}{D}-\frac{1}{(y-z)^{2}}\big]=-\frac{(y-x)^{2}-(y-z)^{2}}{(y-z)^{2}D}$ and $\frac{1}{1-\xi}\big[\frac{1}{D}-\frac{1}{(y-x)^{2}}\big]=\frac{(y-x)^{2}-(y-z)^{2}}{(y-x)^{2}D}$, $D=\xi(y-x)^{2}+(1-\xi)(y-z)^{2}$: with the recoil phase kept in the subtractions the two plus-prescription terms become the regular weights $\frac{(y-x)^{2}}{(y-z)^{2}}B_{2}$, $\frac{(y-z)^{2}}{(y-x)^{2}}B_{1}$ (instead of $B_{2}$, $B_{1}$) plus the two recoil terms. The $q^{+}$ integral of $N_{2a}$ ($\int_{k^{+}}^{0}dq^{+}$, $q^{+}=k^{+}-p^{+}$) is the $B_{1}$ part. The $\frac{11}{6}$ term is the collinear subtraction of $N_{2b\,\rm rest}$, $-\int_{0}^{1}d\xi\,[\xi(1-\xi)-2]=\frac{11}{6}$. Numerically, $\int_{\lambda}^{1-\lambda}d\xi\,[\text{row}]-L\,[-B_{2}/(y-z)^{2}-e^{-is}B_{1}/(y-x)^{2}]-\mathcal{F}$ (contracted with $K^{i}(x'-w')K^{m}(z-x)$) $=3.5\times10^{-4},\,3.5\times10^{-6},\,3.5\times10^{-8}$ for $\lambda=10^{-3},10^{-5},10^{-7}$ (three random points, same rate). Regular term at large distance: $R^{2}\times$(angular mean) $=(-1.2+2.0i)\times10^{-3},\,(-0.1+0.8i)\times10^{-3},\,3\times10^{-4}i,\,1\times10^{-4}i$ for the gluon at $x$ at distance $R=30,100,300,1000$, and $(-1.7+6.5i)\times10^{-3},\,(-0.2+2.0i)\times10^{-3},\,7\times10^{-4}i,\,2\times10^{-4}i$ for the gluon at $z$; also $R^{2}\times$(mean modulus) $\to0$ ($8\times10^{-3}\dots1\times10^{-3}$): no $1/R^{2}$ tail. With the unmodified weights $B_{2}$, $B_{1}$ the same quantity grows ($0.043,\,0.072,\,0.100$ at $R=30,100,300$). Recoil terms at large distance: $R^{2}\times$(angular mean) $=-0.1835,\,-0.2462,\,-0.3031,\,-0.3656$ against $\frac12K(x'-w')\cdot K(y-z)\,[-\gamma_{E}-\log\frac{|k|R}{2}]=-0.1838,\,-0.2462,\,-0.3031,\,-0.3655$ (soft gluon at $x$), and $-0.8915,\,-1.1947,\,-1.4708,\,-1.7743$ against $\frac12K(x'-w')\cdot K(y-x)\,[\dots]=-0.8919,\,-1.1946,\,-1.4709,\,-1.7736$ (soft gluon at $z$). $z\to x$: $K^{i}K^{m}\mathcal{F}^{im}\simeq c(\hat r)/(x-z)^{2}$, and $(x-z)^{2}\times$angular mean of [$K^{i}K^{m}\mathcal{F}^{im}$ $+$ $\frac{11}{6}$ term] $=5.8\times10^{-3},\,5.8\times10^{-5},\,5.8\times10^{-7}$ at $|x-z|=10^{-1},10^{-2},10^{-3}$ (without the $\frac{11}{6}$ term: $0.92$): in $N_{2b\,\rm rest}$ the collinear point is integrable with the angular integral first, as in the checked row; in $N_{2a}$ the colour factor vanishes at $z=x$. $y\to z$, $y\to x$: $1/|\cdot|$.]}

\noindent\textbf{DGLAP and RC.} Group I has no term of DGLAP-evolution form and no running-coupling term: the splitting-function weights $\xi(1-\xi)$, $\frac{1}{\xi}$, $\frac{1}{1-\xi}$ of $N_{2a}$, $N_{2b}$ are integrated over $0<\xi<1$ in $\mathcal{F}^{im}$ (the end-point singularities are the two recoil terms and the JIMWLK terms), the collinear logarithm integrated over $\xi$ is $N_{2b\,\rm UV}$, and it cancels against Row II.

\subsection{Large-$z$ tails of the Group I finite terms}\label{sec:tailsI}
\noindent Terms with a $1/z^{2}$ tail at large position of the soft gluon: the recoil terms $f(0)\,\mathcal{E}(s)$, the second line of \eqref{eq:III}, and (odd in the angle) the first line of \eqref{eq:III}. Table: coefficient $T$ in
\begin{align}
\text{term}\;\xrightarrow[|z_{\rm soft}|\to\infty]{}\;\frac{1}{(2\pi)^{3}}\,\frac{g^{4}}{16\pi^{5}}\,\frac{1}{k^{+}}\;e^{-ik\cdot(w_{\rm hard}-w')}\;\frac{T}{z_{\rm soft}^{2}}\;\mathcal{E}\big(k\cdot(z_{\rm soft}-w_{\rm hard})\big)\,,\notag
\end{align}
split by the degree in the Wilson line at the soft position, $T=T^{(0)}+T^{(1)}$. All terms as written have the recoil on the side of $w$ (amplitude with the $g^{3}$ bracket); their complex conjugates have it on the side of $w'$, with $\mathcal{E}^{*}$. Abbreviations: $P\equiv K(x'-w')$; $\mathcal{L}^{a}\equiv\mathcal{L}^{a}_{x'w'}$; $\Delta U^{ab'}\equiv U^{ab'}(x')-U^{ab'}(w')$; $\{dd\}\equiv f^{abc}U^{cd}(y)U^{be}(y')\{\rho^{d}(y),\rho^{e}(y')\}$.
\begin{center}
\renewcommand{\arraystretch}{1.5}
\begin{tabular}{|l|c|c|l|}
\hline
term & soft & hard & $T^{(0)}$ (no $U$ at the soft position) \\ \hline
III, 3rd line & $z$ & $w$ & $-2i\,P\!\cdot\!K(y-w)\,\mathcal{L}^{a}f^{abd}U^{de}(x)[U^{bc}(w)-U^{bc}(y)]\rho^{e}(x)\rho^{c}(y)$ \\
III, 2nd line $^{(*)}$ & $z$ & $w$ & $-2i\,P\!\cdot\!K(x-w)\,\mathcal{L}^{a}f^{abd}U^{de}(x)U^{bc}(y)\rho^{e}(x)\rho^{c}(y)$ \\
V $N_{1}$, $\mathcal{E}(s)$ & $x$ & $z$ & $+\frac{i}{2}\,P\!\cdot\!K(y-z)\,\mathcal{L}^{a}\{dd\}$ \\
V $N_{1}$, $\mathcal{E}(-s)$ & $z$ & $x$ & $-\frac{i}{2}\,P\!\cdot\!K(y'-x)\,\mathcal{L}^{a}\{dd\}$ \\
V $N_{2a}$, $\mathcal{E}(s)$ & $x$ & $z$ & $0$ \\
V $N_{2a}$, $\mathcal{E}(-s)$ & $z$ & $x$ & $+\frac{N_{c}}{2}\,P\!\cdot\!K(y-x)\,\Delta U^{ab'}U^{af}(x)\rho^{b'}(x')\rho^{f}(y)$ \\
V $N_{2b\,\rm rest}$, $\mathcal{E}(s)$ & $x$ & $z$ & $+\frac{N_{c}}{2}\,P\!\cdot\!K(y-z)\,\Delta U^{ab'}U^{af}(y)\rho^{b'}(x')\rho^{f}(y)$ \\
V $N_{2b\,\rm rest}$, $\mathcal{E}(-s)$ & $z$ & $x$ & $+\frac{N_{c}}{2}\,P\!\cdot\!K(y-x)\,\Delta U^{ab'}[U^{af}(y)-U^{af}(x)]\rho^{b'}(x')\rho^{f}(y)$ \\ \hline
\end{tabular}

\medskip
\begin{tabular}{|l|c|c|l|}
\hline
term & soft & hard & $T^{(1)}$ (linear in $U$ at the soft position) \\ \hline
III, 3rd line & $z$ & $w$ & $0$ \\
III, 2nd line $^{(*)}$ & $z$ & $w$ & $+2i\,P\!\cdot\!K(x-w)\,\mathcal{L}^{a}f^{abd}U^{de}(x)U^{bc}(z)\rho^{e}(x)\rho^{c}(y)$ \\
V $N_{1}$, $\mathcal{E}(s)$ & $x$ & $z$ & $-\frac{i}{2}\,P\!\cdot\!K(y-z)\,\mathcal{L}^{a}f^{adc}U^{db}(x)U^{ce}(z)\{\rho^{e}(y),\rho^{b}(y')\}$ \\
V $N_{1}$, $\mathcal{E}(-s)$ & $z$ & $x$ & $+\frac{i}{2}\,P\!\cdot\!K(y'-x)\,\mathcal{L}^{a}f^{adc}U^{db}(x)U^{ce}(z)\{\rho^{e}(y),\rho^{b}(y')\}$ \\
V $N_{2a}$, $\mathcal{E}(s)$ & $x$ & $z$ & $+\frac12\,P\!\cdot\!K(y-z)\,\Delta U^{ab'}[-f^{adc}f^{fbe}U^{db}(x)U^{ce}(z)+N_{c}U^{af}(x)]\rho^{b'}(x')\rho^{f}(y)$ \\
V $N_{2a}$, $\mathcal{E}(-s)$ & $z$ & $x$ & $-\frac12\,P\!\cdot\!K(y-x)\,\Delta U^{ab'}f^{adc}f^{fbe}U^{db}(x)U^{ce}(z)\rho^{b'}(x')\rho^{f}(y)$ \\
V $N_{2b\,\rm rest}$, $\mathcal{E}(s)$ & $x$ & $z$ & $-\frac{N_{c}}{2}\,P\!\cdot\!K(y-z)\,\Delta U^{ab'}U^{af}(x)\rho^{b'}(x')\rho^{f}(y)$ \\
V $N_{2b\,\rm rest}$, $\mathcal{E}(-s)$ & $z$ & $x$ & $0$ \\ \hline
\end{tabular}
\end{center}
$^{(*)}$ For the second line of \eqref{eq:III} the function is $-\log\frac{|z-w|}{|x-w|}$ instead of $\mathcal{E}$. Under the $w$ integral and after the angular average the two are equivalent: $\langle\mathcal{E}(k\cdot z)\rangle_{\rm angle}\to-\gamma_{E}-\log\frac{|k||z|}{2}$ and
\begin{align}
\int_{w}e^{-ik\cdot w}\,K^{i}(x-w)\Big[\gamma_{E}+\log\frac{|k|\,|x-w|}{2}\Big]=0\qquad\Big(\int_{0}^{\infty}dt\,J_{1}(t)\Big[\gamma_{E}+\log\frac{t}{2}\Big]=0\Big)\,.\label{eq:J1}
\end{align}
The first line of \eqref{eq:III} has the tail \eqref{eq:saw} (zero angular mean; {\color{magenta}OPEN I-1}), not of this form. Row II, Row IV and the regular terms of Row V have no $1/z^{2}$ tail.
{\color{gray}\small [Tails: analytic from the kernels ($K(x,w,z)\to\frac{1}{z^{2}}$, $K(y,w,z)\to\frac{1}{z^{2}}$, $K(z-x)\cdot K(y'-x)\to\frac{1}{x^{2}}$, $K(z-x)\cdot K(y-z)\to-\frac{1}{z^{2}}$, $K^{m}(z-x)B_{2}^{im}/(y-z)^{2}\to-\frac{K^{i}(y-z)}{2x^{2}}$, $K^{m}(z-x)B_{1}^{im}/(y-x)^{2}\to-\frac{K^{i}(y-x)}{2z^{2}}$); checked numerically for $N_{2a}$, $N_{2b}$ (numbers under Eq.~\eqref{eq:V2b}) and for the second line of \eqref{eq:III}. Eq.~\eqref{eq:J1}: numerically $-1.4\times10^{-8}$.]}

\noindent\textbf{Sums of the tails} (numerical, with the matrix engine of the note on the flagged terms: explicit $SU(2)$ and $SU(3)$ matrices, random Wilson lines, ordered products reduced to the symmetrised two-$\rho$ part; maximum norm; the weights of the table).
\begin{itemize}
\item $T^{(0)}$, recoil on $w$: the entries of Row V ($N_{1}$, $N_{2a}$, $N_{2b\,\rm rest}$) and of the second line of Row III cancel among themselves: $|{\rm sum}|=5\times10^{-8}$ ($SU(2)$; single entries $0.46$--$1.74$), $1\times10^{-8}$ ($SU(3)$; $0.80$--$2.02$). By \eqref{eq:J1} this is a complete cancellation of the $1/z^{2}$ tail if the $w$ integral and the angular integral are done before the radial $z$ integral. The entry of the third line of Row III ($2.86$ for $SU(2)$, $1.60$ for $SU(3)$) is not cancelled inside Group I; it is cancelled by Group III Row V on the same side ($|{\rm sum}|=5\times10^{-8}$, $1\times10^{-8}$). The same holds for the c.c. terms (recoil on $w'$): $6\times10^{-9}$, $1\times10^{-8}$.
\item $T^{(1)}$ (linear in $U$ at the soft position): not cancelled inside Group I (sum $3.12$ for $SU(2)$, $3.19$ for $SU(3)$), and not together with Group III Rows III, V, VI on the same side ($5.22$, $2.37$). It cancels only between the two sides, Group I $+$ Group III with the c.c. terms: $2\times10^{-8}$, $1\times10^{-8}$.
\end{itemize}
\noindent{\color{magenta}\fbox{OPEN I-2} (a) The recoil terms of Rows III and V and the second line of \eqref{eq:III} are not integrable at large soft position one by one ($\frac{T}{z^{2}}\log|z|$); they have to be added under one integral over the soft position, with the hard position as the common variable $w$ (table). Inside Group I this makes the $T^{(0)}$ part finite except for the third line of Row III, which needs Group III Row V; whoever assembles the file has to check that Group III Row V is written with the same function $\mathcal{E}$ of the same argument. (b) For the part linear in $U$ at the soft position the cancellation is between the $w$ side and the $w'$ side, where the functions are $\mathcal{E}$ and $\mathcal{E}^{*}$: what is left is $\frac{1}{z^{2}}\big[\mathcal{E}-\mathcal{E}^{*}\big]T^{(1)}\to-\frac{i\pi}{z^{2}}\,{\rm sgn}(k\cdot z)\,T^{(1)}(z)$ ($T^{(1)}$: sum of the $w$-side entries of Groups I and III). This is integrable only if $U(z)$ has a direction-independent limit at large $|z|$; not decided here. (c) The term that remains at large distance after all cancellations is the non-singlet term $L\int\frac{d^{2}z}{z^{2}}$ of the JIMWLK formula; it is not part of the finite terms.}

\subsection{What was not checked in this section}
\noindent (1) No complete transverse integration with model Wilson lines was done; the checks are pointwise limits at random points, integrals over the soft position with a test function, and scaling exponents at the coincidence points and at large distance. (2) Row II, Eq.~\eqref{eq:IIb}: the kernel is given in momentum space; the convergence at large $p$, $q$ rests on the phases; the coordinate-space kernel was not derived. (3) The sums of the tails use the $\log\frac{k^{+}}{\Lambda}$ coefficients of the checked rows (two-$\rho$ sector, split by side and by degree in $U$); in the three-$\rho$ sector only the total over Group I and Group III was checked earlier, not the split. (4) The $1/z^{2}$ tails of the Row V $N_{1}$ recoil terms are analytic only.

\section{Group II}
\noindent\textbf{Finite terms of Groups II and III at $\Lambda\to0$, $\vee\to\infty$.} What is listed here is, row by row, [checked row with the cut-offs $\Lambda<p^{+}<\vee-k^{+}$] $-$ $\log\frac{\vee}{k^{+}}\times$[BFKL-list term of the row] $-$ $\big[\log\frac{k^{+}}{\Lambda}+\log\big(1-\frac{k^{+}}{\vee}\big)\big]\times$[JIMWLK-list term of the row], in the limit; the terms with $\log\big(1-\frac{k^{+}}{\vee}\big)$ go with the $\log\frac{k^{+}}{\Lambda}$ term of the same row and are not listed. Conventions: phase $e^{-ik\cdot(w-w')}$ and descendants; c.c.\ is implied for Group III, not for Group II; $\rho^{a}(x)$ are operators and keep the order of the checked rows; no $\frac{1}{\epsilon}$ is written; logarithms are $\log|\dots|$. Shorthand used in both groups:
\begin{align}
\mathcal{K}^{i}(r)\equiv\frac{r^{i}}{r^{2}}\,,\qquad
\mathcal{T}^{ij}(x,z,w)\equiv\frac{\delta^{ij}}{2\,(z-w)^{2}}+\mathcal{K}^{i}(z-w)\,\mathcal{K}^{j}(x-z)-\mathcal{K}^{i}(x-w)\,\mathcal{K}^{j}(x-z)-\mathcal{K}^{i}(x-w)\,\mathcal{K}^{j}(z-w)\,.\notag
\end{align}
{\color{gray}\small Grey lines: origin of the equation and what was checked numerically (scripts in \texttt{finite/g23}). ``Checked row'' = \texttt{Group\_II\_checked}, \texttt{Group\_II\_Row\_IV\_checked}, \texttt{Group\_III\_checked}; ``limits file'' = \texttt{Other\_terms\_limits}.}\par
{\color{magenta}Magenta: what is not finite or not settled (OPEN II/III-1 to II/III-4).}
\subsection{Row I}
\subsubsection{First part}
\begin{align}
 \frac{d^{3}N_{1a}}{d^{3}k}  & =  -\,\frac{1}{(2\pi)^{3}}\,\frac{ig^{4}}{16\pi^{5}}\,\frac{1}{k^{+}}\int_{x,y,z,w,x',w^{\prime}} e^{-ik\cdot(w-w^{\prime})}\,\mathcal{K}^{i}(x'-w')\,\mathcal{K}^{j}(y-z)\,\log\left[\frac{(x-z)^{2}}{(x-w)^{2}}\right]\,\mathcal{T}^{ij}(x,z,w)\notag \\
& \times \Big[U^{ac'}(w')\,\{\rho^{c'}(x'),\rho^{b}(y)\}-U^{c'b}(z)\,U^{ad'}(x')\,U^{c'e'}(y)\, \{\rho^{d'}(x'),\rho^{e'}(y)\}\Big]\Big[f^{cbd}\,U^{ac}(w)\,\rho^{d}(x)-f^{acd}\, U^{cb}(z)\,U^{de}(x)\,\rho^{e}(x)\Big]
\end{align}
\begin{align}
\frac{d^{3}N_{1b}}{d^{3}k} & =  +\,\frac{1}{(2\pi)^{3}}\,\frac{ig^{4}}{16\pi^{5}}\,\frac{1}{k^{+}}\int_{x,y,z,w,x',w^{\prime}} e^{-ik\cdot(w-w^{\prime})}\,\mathcal{K}^{i}(x'-w)\,\mathcal{K}^{j}(y-z)\,\log\left[\frac{(x-z)^{2}}{(x-w')^{2}}\right]\,\mathcal{T}^{ij}(x,z,w')\notag \\
& \times \Big[f^{c'bd'}\,U^{ac'}(w')\,\rho^{d'}(x)-f^{ac'd'}\, U^{c'b}(z)\,U^{d'e'}(x)\,\rho^{e'}(x)\Big]\Big[U^{ac}(w)\,\{\rho^{c}(x'),\rho^{b}(y)\}-U^{cb}(z)\,U^{ad}(x')\,U^{ce}(y)\, \{\rho^{d}(x'),\rho^{e}(y)\}\Big]
\end{align}
{\color{gray}\small From the first part of Row I of the checked row (the $p^{+}$ integral before it is done); same as the limits file. Checked at a random point: [$p^{+}$ integral from $\Lambda$ to $\vee-k^{+}$] $-\log\frac{\vee}{k^{+}}\big[\mathcal{K}^{i}(z-w)-\frac12\mathcal{K}^{i}(x-w)\big]\mathcal{K}^{j}(x-z)$ $-\log\frac{k^{+}}{\Lambda}\,\mathcal{K}^{i}(x-w)\big[\mathcal{K}^{j}(z-w)+\frac12\mathcal{K}^{j}(x-z)\big]$ minus the kernel above $=9.0\times10^{-4}$, $9.1\times10^{-6}$, $9.1\times10^{-8}$, $9.1\times10^{-10}$ for $\frac{\Lambda}{k^{+}}=\frac{k^{+}}{\vee}=10^{-2},10^{-4},10^{-6},10^{-8}$ (value $-0.0407$). $N_{1b}$ is $N_{1a}$ with $w\leftrightarrow w'$ in the kernel and the opposite sign.}\par
\noindent Integrability (each of $N_{1a}$, $N_{1b}$ by itself): $x\to z$ and $x\to w$: logarithmic only (in $\mathcal{T}^{ij}$ the terms with $\mathcal{K}^{j}(x-z)$, resp. with $\mathcal{K}^{i}(x-w)$, cancel at leading order; numerically the kernel grows like $r^{-0.09}$ between $r=10^{-4}$ and $10^{-6}$); $z\to w$: $1/|z-w|$ (the logarithm vanishes linearly); $y\to z$: $1/|y-z|$; large $z$: $\log z^{2}/|z|^{3}$ (no $1/z^{2}$ tail; numerically $|z|^{-2.9}$); large $w$: $\log w^{2}/|w|$ times the phase, as at leading order.
\subsubsection{Second part}
\noindent Notation for this part: $a\equiv x-w$, $a'\equiv x'-w'$, $b\equiv x-z$, $b'\equiv x'-z$, $c\equiv z-w$, $c'\equiv z-w'$; $A\equiv a^{2}$, $A'\equiv a'^{2}$, $B\equiv b^{2}$, $B'\equiv b'^{2}$, $C\equiv c^{2}$, $C'\equiv c'^{2}$; $\mathcal{D}\equiv A'B-AB'$; $\ell\equiv\log\frac{(x-z)^{2}}{(x-w)^{2}}$, $\ell'\equiv\log\frac{(x'-z)^{2}}{(x'-w')^{2}}$.
\begin{align}
\frac{d^{3}N_{2}}{d^{3}k} & = \frac{1}{(2\pi)^{3}}\, \frac{g^{4}}{8\pi^{5}}\,\frac{1}{k^{+}}\,\int_{z,w,w^{\prime},x,x'}\, e^{-ik\cdot(w-w^{\prime})}\,\Big[\mathfrak{K}_{2a}+\mathfrak{K}_{2b}+\mathfrak{K}_{2c}+\mathfrak{K}_{2d}+\mathfrak{K}_{2e}+\mathfrak{K}_{2f}\Big]\notag \\
& \times\Big[f^{c'bd'}\,f^{cbd}\, U^{ac'}(w')\,U^{ac}(w)\,\rho^{d'}(x')\rho^{d}(x)- f^{c'bd'}\,f^{acd}\,U^{ac'}(w')\,U^{cb}(z)\,U^{de}(x)\,\rho^{d'}(x')\rho^{e}(x) \notag \\
& \quad -f^{ac'd'}\,f^{cbd}\,U^{ac}(w)\,U^{c'b}(z)\,U^{d'e'}(x')\,\rho^{e'}(x')\,\rho^{d}(x) +N_{c}\,U^{de'}(x')\,U^{de}(x)\,\rho^{e'}(x')\,\rho^{e}(x)\Big]
\end{align}
\begin{align}
\mathfrak{K}_{2a} & =\frac{\ell-\ell'}{\mathcal{D}}\,\Bigg[\frac{a\cdot c'\;b'\cdot c}{C\,C'}-\frac{a\cdot a'\;c\cdot b'}{2\,C\,A'}+\frac{b\cdot b'\;a\cdot c'}{2\,B\,C'}-\frac{b\cdot b'\;a\cdot a'}{4\,B\,A'}
+\frac{a'\cdot c\;b\cdot c'}{C\,C'}+\frac{a'\cdot c\;b\cdot b'}{2\,C\,B'}-\frac{b\cdot c'\;a\cdot a'}{2\,A\,C'}-\frac{a\cdot a'\;b\cdot b'}{4\,A\,B'}\Bigg]\\
\mathfrak{K}_{2b} & =\frac{1}{\mathcal{D}}\left[\frac{B'}{A'}\,\ell'-\frac{B}{A}\,\ell\right]\Bigg[\frac{a\cdot a'\;c\cdot c'}{C\,C'}+\frac{a\cdot a'\;b'\cdot c}{2\,C\,B'}+\frac{a\cdot a'\;b\cdot c'}{2\,B\,C'}+\frac{b\cdot b'\;a\cdot a'}{4\,B\,B'}\Bigg]\\
\mathfrak{K}_{2c} & =\frac{1}{\mathcal{D}}\left[\frac{A'}{B'}\,\ell'-\frac{A}{B}\,\ell\right]\Bigg[\frac{b\cdot b'\;c\cdot c'}{C\,C'}-\frac{b\cdot b'\;a'\cdot c}{2\,C\,A'}-\frac{a\cdot c'\;b\cdot b'}{2\,A\,C'}+\frac{b\cdot b'\;a\cdot a'}{4\,A\,A'}\Bigg]\\
\mathfrak{K}_{2d} & =\frac{1}{\mathcal{D}}\left[\frac{B\,\ell}{B-A}+\frac{B'\,\ell'}{A'-B'}\right]\Bigg[\frac{a\cdot c\,(B'-A')}{2\,C\,C'}+\frac{a\cdot b\,(B'-A')}{4\,B\,C'}+\frac{a'\cdot c'\,(B-A)}{2\,C\,C'}+\frac{a'\cdot b'\,(B-A)}{4\,C\,B'}\Bigg]\\
\mathfrak{K}_{2e} & =-\frac{1}{\mathcal{D}}\left[\frac{A\,\ell}{B-A}+\frac{A'\,\ell'}{A'-B'}\right]\Bigg[\frac{b\cdot c\,(B'-A')}{2\,C\,C'}-\frac{a\cdot b\,(B'-A')}{4\,A\,C'}+\frac{b'\cdot c'\,(B-A)}{2\,C\,C'}-\frac{a'\cdot b'\,(B-A)}{4\,C\,A'}\Bigg]\\
\mathfrak{K}_{2f} & =\frac{d_{\perp}}{4\,C\,C'}\left[-\frac{A\,B\,(B'-A')}{(B-A)\,\mathcal{D}}\,\ell\;-\;1\;+\;\frac{A'\,B'\,(B-A)}{(B'-A')\,\mathcal{D}}\,\ell'\right],\qquad d_{\perp}=2
\end{align}
{\color{gray}\small From the second part of Row I of the checked row ($N_{2a},\dots,N_{2f}$ of the limits file, written with one colour factor; the first term of $N_{2b}$, $\propto\log(1-\frac{k^{+}}{\vee})$, has gone with the JIMWLK term; the $-1$ in $\mathfrak{K}_{2f}$ is the power term $k^{+}\big[\frac{1}{\vee}-\frac{1}{k^{+}}\big]\to-1$). Checked at a random point: $k^{+}\int_{\Lambda}^{\vee-k^{+}}dp^{+}$[integrand of the checked row] $-\log\frac{\vee}{k^{+}}\,\frac{1}{BB'}[\text{bracket of }\mathfrak{K}_{2c}]-\log\frac{k^{+}}{\Lambda}\,\frac{1}{AA'}[\text{bracket of }\mathfrak{K}_{2b}]$ minus $\sum\mathfrak{K}_{2\cdot}$ $=4.8\times10^{-3}$, $6.5\times10^{-5}$, $6.5\times10^{-7}$, $6.5\times10^{-9}$ for $\frac{\Lambda}{k^{+}}=\frac{k^{+}}{\vee}=10^{-2},10^{-4},10^{-6},10^{-8}$ (value $-0.1798$; single kernels up to $1.3$).}\par
\noindent\textbf{What has to stay under one integral.} (1) The two terms ($\propto\ell$ and $\propto\ell'$) of each of $\mathfrak{K}_{2b},\dots,\mathfrak{K}_{2f}$: on the surface $\mathcal{D}=0$ each term has a simple pole and the residues cancel within each $\mathfrak{K}$ (in $\mathfrak{K}_{2a}$ the factor $\ell-\ell'$ vanishes there); on $B=A$ and on $B'=A'$ the logarithms vanish. Each $\mathfrak{K}_{2\cdot}$ is then finite on these surfaces [numerically: at distance $10^{-3}\dots10^{-7}$ from $\mathcal{D}=0$ the single terms grow like $1/\mathcal{D}$ (up to $10^{5}$), each $\mathfrak{K}$ is constant to four digits]. (2) The three terms of $\mathfrak{K}_{2f}$: at $z\to w$ ($z\to w'$) the first (third) term tends to $+\frac{d_{\perp}}{4CC'}$ and cancels the $-1$; $\mathfrak{K}_{2f}={\cal O}(1/|z-w|)$, ${\cal O}(1/|z-w'|)$, so no pole multiplies $d_{\perp}$ and $d_{\perp}=2$. (3) \emph{All six kernels have to be integrated over $z$ together:} at large $z$, $\mathfrak{K}_{2a}\to\frac{\rm const}{z^{2}}$, $\mathfrak{K}_{2b},\mathfrak{K}_{2c}\to\pm\frac{\log z^{2}}{z^{2}}$, while $\sum\mathfrak{K}_{2\cdot}\sim\log z^{2}/|z|^{3}$ [$R^{2}\sum\mathfrak{K}=-2.7\times10^{-2}$, $-5.8\times10^{-4}$, $-8.9\times10^{-6}$ at $|z|=R=10^{2},10^{4},10^{6}$; $R^{2}\mathfrak{K}_{2b}=0.49,\,1.00,\,1.49$]. With (1)--(3) the remaining singularities are integrable: $x\to w$, $x\to z$, $x'\to w'$, $x'\to z$, $z\to w$, $z\to w'$: $1/r$ (times a logarithm); $z\to w=w'$: $1/r^{3.3}$ in four dimensions ($z$ and $w-w'$), integrable; large $w$: $1/|w|$ times the phase, as at leading order.
\subsection{Row II}
\textsc{Only 4 $\rho$'s}
\subsection{Row III}
\textsc{Only 4 $\rho$'s}
\subsection{Row IV}
\noindent Notation: $\mathcal{C}\equiv\frac{1}{(2\pi)^{3}}\frac{g^{2}}{2\pi^{2}}\frac{g^{2}}{8\pi^{2}}\frac{1}{k^{+}}=\frac{1}{(2\pi)^{3}}\frac{g^{4}}{16\pi^{4}}\frac{1}{k^{+}}$,
\begin{align}
\mathcal{F}(\xi) & \equiv\frac{1}{\xi^{2}}\int_{x,y,x',y'}e^{-ik\cdot(y-y')\frac{1}{\xi}}\Big[2\gamma_{E}+\log\frac{(y'-y)^{2}\mu^{2}_{\overline{MS}}}{4\,\xi^{2}}\Big]\,\mathcal{K}(x-y)\cdot\mathcal{K}(x'-y')\notag\\ & \qquad\times\left[U^{ce'}(y') - U^{ce'}(x')\right]\left[U^{ce}(y) - U^{ce}(x)\right]\rho^{e'}(x')\rho^{e}(x)\,,\notag\\
\mathcal{X} & \equiv\int_{x,x',z}\mathcal{K}(x-z)\cdot\mathcal{K}(x'-z)\left[U^{ce'}(z) - U^{ce'}(x')\right]\left[U^{ce}(z) - U^{ce}(x)\right]\rho^{e'}(x')\rho^{e}(x)\,,\qquad \mathcal{F}(0)\equiv\lim_{\xi\to0}\mathcal{F}(\xi)=-\frac{4\pi}{k^{2}}\,\mathcal{X}\,.\notag
\end{align}
In $\mathcal{F}(\xi)$ the integral over $y-y'$ is done before the integral over $\xi$. The BFKL term of this row is $+2N_{c}\,\mathcal{C}\,\frac{4\pi}{k^{2}}\log\frac{\vee}{k^{+}}\,\mathcal{X}$ (a pure $\log\frac{\vee}{k^{+}}$), its JIMWLK-type term is $-2N_{c}\,\mathcal{C}\big[\log\frac{k^{+}}{\Lambda}+\log(1-\frac{k^{+}}{\vee})\big]\mathcal{F}(1)$.
\subsubsection{(i) DGLAP term}
\begin{align}
\frac{d^{3}N_{\rm DGLAP}}{d^{3}k} & =-\mathcal{C}\int_{0}^{1}d\xi\;\widetilde{\mathcal P}_{gg}(\xi)\,\mathcal{F}(\xi)\,,\qquad
\widetilde{\mathcal{P}}_{gg}(\xi) = 2\,N_{c}\left[\xi(1-\xi) + \frac{\xi}{(1-\xi)_{+}} + \Big[\frac{1-\xi}{\xi}\Big]_{+}\right]+\frac{\beta_{0}}{2}\,\delta(1-\xi)\,,\notag\\
& \equiv-2N_{c}\,\mathcal{C}\int_{0}^{1}d\xi\,\Bigg\{\xi(1-\xi)\,\mathcal{F}(\xi)+\frac{\xi\,\mathcal{F}(\xi)-\mathcal{F}(1)}{1-\xi}+\frac{(1-\xi)\,\mathcal{F}(\xi)-\mathcal{F}(0)}{\xi}\Bigg\}\;-\;\mathcal{C}\,\frac{\beta_{0}}{2}\,\mathcal{F}(1)
\end{align}
Both plus prescriptions act on the whole function of $\xi$ that multiplies $\frac{1}{1-\xi}$ resp. $\frac{1}{\xi}$, i.e. on $\xi\mathcal{F}(\xi)$ and on $(1-\xi)\mathcal{F}(\xi)$; the subtraction at $\xi=0$ is the compensating term $+\frac{2N_{c}}{\xi}\frac{4\pi}{k^{2}}\mathcal{X}$. With this reading there is no remnant term: the term $-\big(1-\frac{k^{+}}{\vee}\big)\times$(BFKL structure) of the limits file is contained in it. [If $\big[\frac{1-\xi}{\xi}\big]_{+}$ is read as $\frac{1-\xi}{\xi}\big[\mathcal{F}(\xi)-\mathcal{F}(0)\big]$, as in the note in \texttt{Divergent\_terms\_checked}, the term $-2N_{c}\,\mathcal{C}\,\frac{4\pi}{k^{2}}\mathcal{X}$ has to be added; the two forms are identical.]\par
{\color{gray}\small From the checked Row IV file (equation after ``Relabelling coordinates'', bracket without $\frac{2}{\epsilon}$) and \texttt{Cancellation\_of\_flagged\_terms} (Sec.~2.3). $\xi\to0$: $\mathcal{F}(\xi)-\mathcal{F}(0)={\cal O}(\xi)$, and ${\cal O}(\xi^{2})$ in the part symmetric under $y\leftrightarrow y'$; the lower limit is $0$ [test function with an odd part: $|\mathcal{F}(\xi)-\mathcal{F}(0)|=0.44$, $0.086$, $0.027$ at $\xi=0.1$, $0.03$, $0.01$; real (even) part $0.33$, $0.028$, $0.0031$]. At fixed $y-y'$ the $\xi$ integral does not exist. Numerical check with (i) and (iii) together, below.}
\subsubsection{(ii) Running-coupling term}
\begin{align}
\frac{d^{3}N_{\rm RC}}{d^{3}k} & = +\,\mathcal{C}\,\frac{\beta_{0}}{2}\,\mathcal{F}(1)= +\frac{1}{(2\pi)^{3}}\frac{g^{2}}{2\pi^{2}}\frac{g^{2}\beta_{0}}{16\pi^{2}}\frac{1}{k^{+}}\int_{x,y,x',y'}e^{-ik\cdot(y-y')}\Big[2\gamma_{E}+\log\frac{(y'-y)^{2}\mu^{2}_{\overline{MS}}}{4}\Big]\mathcal{K}(x-y)\cdot\mathcal{K}(x'-y')\notag \\ & \qquad\times\left[U^{ce'}(y') - U^{ce'}(x')\right]\left[U^{ce}(y) - U^{ce}(x)\right]\rho^{e'}(x')\rho^{e}(x)
\end{align}
It is minus the $\delta(1-\xi)$ term of $N_{\rm DGLAP}$ (added and subtracted, one scale $\mu$): $N_{\rm DGLAP}+N_{\rm RC}$ does not depend on $\beta_{0}$. Combined with the leading order ($\alpha_{s}=\frac{g^{2}}{4\pi}$, $r\equiv y'-y$):
\begin{align}
\frac{d^{3}N_{\rm LO}}{d^{3}k}+\frac{d^{3}N_{\rm RC}}{d^{3}k}&=\frac{1}{(2\pi)^{3}}\,\frac{g^{2}(\mu)}{2\pi^{2}}\,\frac{1}{k^{+}}\int_{x,y,x',y'}e^{-ik\cdot(y-y')}\left[1+\frac{\alpha_{s}(\mu)\,\beta_{0}}{4\pi}\log\frac{r^{2}\,\mu^{2}_{\overline{MS}}\,e^{2\gamma_{E}}}{4}\right]\mathcal{K}(x-y)\cdot\mathcal{K}(x'-y')\notag\\
&\qquad\times\left[U^{ce'}(y')-U^{ce'}(x')\right]\left[U^{ce}(y)-U^{ce}(x)\right]\rho^{e'}(x')\,\rho^{e}(x)\,,
\end{align}
and to this order the square bracket is $\alpha_{s}(Q^{2})/\alpha_{s}(\mu^{2})$ with $Q^{2}=\frac{4\,e^{-2\gamma_{E}}}{r^{2}}$: the coupling of the leading-order term runs at the scale $4e^{-2\gamma_{E}}/(y'-y)^{2}$. This holds only together with the $\delta(1-\xi)$ term kept in $\widetilde{\mathcal P}_{gg}$ at the same $\mu$.\par
{\color{gray}\small From the checked Row IV file and \texttt{Divergent\_terms\_checked} (Sec.~``Running coupling''); no dependence on $\Lambda$, $\vee$.}
\subsubsection{(iii) Finite term}
\begin{align}
\frac{d^{3}N_{\rm Finite}}{d^{3}k}= &  +\frac{1}{(2\pi)^{3}}\,\frac{g^{4}\,N_{c}}{8\pi^{4}}\,\frac{1}{k^{+}}\int_{0}^{1}d\xi\,\frac{1-\xi}{\xi}\int_{x,y,x',y'}e^{-ik\cdot(y-y')\frac{1}{\xi}}\notag\\
&\times\bigg[\mathcal{K}(x-y)\cdot\mathcal{K}(x'-y')-2\,\frac{\mathcal{K}(x-y)\cdot(y'-y)\;\mathcal{K}(x'-y')\cdot(y'-y)}{(y'-y)^{2}}\bigg]\notag\\
&\times\left[U^{ce'}(y') - U^{ce'}(x')\right]\left[U^{ce}(y) - U^{ce}(x)\right]\rho^{e'}(x')\rho^{e}(x)
\end{align}
{\color{gray}\small From the checked Row IV file (last term); lower limit $0$: after the $y-y'$ integral the integrand is ${\cal O}(\xi)$ (part cut off: ${\cal O}(k^{+2}/\vee^{2})$). At fixed coordinates the $\xi$ integral converges conditionally, with an integrable $-\log|k\cdot(y-y')|$ on the line $k\cdot(y-y')=0$.}\par
\noindent\textbf{Numerical check of (i)--(iii).} With test functions in place of the Wilson-line structure after the integration over $\frac{y+y'}{2}$ [trace part: two Gaussians in $y'-y$ times $(1+c\cdot(y'-y))$; tensor part $\propto\cos2\varphi$ times a Gaussian], in units of $-2N_{c}\mathcal{C}$: (i) without the $\beta_{0}$ term $=13.982+2.096\,i$, (iii) $=-1.839$, $\mathcal{F}(1)=0.218-0.560\,i$, $\mathcal{F}(0)=-7.113$; [row with cut-offs, pole dropped] $-\log\frac{\vee}{k^{+}}\,\mathcal{F}(0)$ $-\big[\log\frac{k^{+}}{\Lambda}+\log(1-\frac{k^{+}}{\vee})\big]\mathcal{F}(1)$ $-$ (i) $-$ (iii) $=0.145$, $1.5\times10^{-2}$, $1.5\times10^{-3}$, $1.5\times10^{-4}$ in modulus for $\frac{\Lambda}{k^{+}}=\frac{k^{+}}{\vee}=10^{-2},\dots,10^{-5}$, i.e. ${\cal O}(\frac{\Lambda}{k^{+}},\frac{k^{+}}{\vee})$. (The $\beta_{0}$ terms of (i) and (ii) cancel identically.) Dropped: the pole $\frac{2}{\epsilon}$ in each bracket $[2\gamma_{E}+\log\dots]$ (collinear, large-distance origin); the pole of the DGLAP term is not cancelled by another term.\par
\noindent Integrability of (i)--(iii): $y'\to y$: logarithm; $x\to y$, $x'\to y'$: $\mathcal{K}(x-y)[U(y)-U(x)]$ is bounded; $\xi\to1$: regular after the subtraction; $\xi\to0$: see above; large distance: not integrable, OPEN II/III-1.\par\smallskip
\noindent{\color{magenta}\fbox{OPEN II/III-1} \textbf{Large-distance tail of Row IV.} With $G_{\infty}(Y)\equiv\int_{x,x'}[U(Y)-U(x')]^{ce'}[U(Y)-U(x)]^{ce}\rho^{e'}(x')\rho^{e}(x)$ and $\varphi$ the angle between $k$ and $Y=\frac{y+y'}{2}$, the terms (i)+(ii)+(iii) behave at large $Y$ like
\begin{align}
{\color{magenta}2N_{c}\,\mathcal{C}\,\frac{4\pi}{k^{2}}\int\frac{d^{2}Y}{Y^{2}}\,\Big[-\frac{11}{6}+\frac16\cos 2\varphi\Big]\,G_{\infty}(Y)}\notag
\end{align}
[(i): $-\frac{11}{12}$, with $\beta_{0}=\frac{11N_{c}}{3}$; (ii): $-\frac{11}{12}$; (iii): $+\frac16\cos2\varphi$]: a logarithmic divergence, $-\frac{11}{6}$ times the large-distance tail of the BFKL structure $2N_{c}\mathcal{C}\frac{4\pi}{k^{2}}\mathcal{X}$, plus a piece that integrates to zero only if $G_{\infty}$ does not depend on the direction of $Y$. $G_{\infty}\neq0$ also for a neutral projectile and in the singlet projection. No term of Groups II, III cancels it; whether it belongs to the large-distance term of the BFKL formula or to the collinear pole that is not written is not decided here. The two coefficients are now confirmed numerically in a model (one point charge at the origin, $[U(y)-U(x)]\to g(y)$ with $g=0$ at the charge, $g=1$ for $\rho\ll|y|\ll R$ and a Gaussian cut-off at $|y|\sim R$, isotropic or elliptic with axes $R$, $3R$; $k\rho=\mu\rho=1$; evaluation in momentum space): [(i)+(ii)+(iii)]$/(2N_{c}\mathcal{C})$ $+\frac{4\pi}{k^{2}}\int\frac{d^{2}Y}{Y^{2}}\big[\frac{11}{6}-\frac16\cos2\varphi\big]g^{2}(Y)$ $=73.75$, $68.34$, $67.09$, $66.78$, $66.70$ for $R/\rho=8,16,32,64,128$ (isotropic; differences decrease like $1/R^{2}$, while the subtracted term goes from $204$ to $602$); elliptic cut-off, $k$ along the short axis, at $45^{\circ}$, along the long axis: $66.43$, $66.67$, $66.91$ at $R/\rho=64$, the spread ($1.40$, $0.66$, $0.24$ at $R/\rho=16,32,64$) decreasing like $\log R/R^{2}$; (iii) alone minus its $\frac16\cos2\varphi$ term: $-0.730$, $-0.728$, $-0.730$, while (iii) itself varies between $-4.02$ and $+2.56$. The statement is for the non-oscillating part: terms that oscillate like $e^{2ik\cdot Y/\xi}$ are of the same order pointwise and integrate to a finite result.}
\section{Group III\, (with c.c.)}
\noindent Functions of $s\equiv k\cdot(z-y')$ used in Rows III, V, VI ($E_{1}$: exponential integral; $t=p^{+}/k^{+}$):
\begin{align}
\widehat{\mathcal{E}}(s) & \equiv\int_{0}^{\infty}\frac{dt}{1+t}\,e^{-its}=e^{is}E_{1}(is)\,,\qquad
\Psi(s)\equiv\int_{0}^{\infty}\frac{dt}{t\,(1+t)}\,\big(e^{-its}-1\big)=-\gamma_{E}-\log|s|-\frac{i\pi}{2}\,{\rm sgn}(s)-\widehat{\mathcal{E}}(s)\,.\notag
\end{align}
$\widehat{\mathcal{E}}(s)\simeq-\gamma_{E}-\log|s|-\frac{i\pi}{2}{\rm sgn}(s)$ at $s\to0$ and $\widehat{\mathcal{E}}(s)\simeq\frac{1}{is}$ at $|s|\to\infty$; $\Psi(s)={\cal O}(s\log|s|)$ at $s\to0$ and $\Psi(s)\simeq-\gamma_{E}-\log|s|-\frac{i\pi}{2}{\rm sgn}(s)$ at $|s|\to\infty$.

\noindent\textbf{How Rows III, V, VI are written.} BFKL side ($p^{+}\to\infty$): the author's split is kept, kernel $\big[\frac{k^{+}}{p^{+}(p^{+}+k^{+})}\big]_{+}+\frac{1}{p^{+}+k^{+}}$ and the term with the constant kernel subtracted at $y'\to z$; the subtracted terms are written with the upper limit $\infty$, which is allowed if the integral over $y'$ is done before the integral over $p^{+}$ (equivalently in the distributional form given). Soft side ($p^{+}\to0$): the subtraction of the plus prescription keeps the recoil phase,
\begin{align}
\int_{\Lambda}^{\vee-k^{+}}\!\!dp^{+}\,\frac{k^{+}\,e^{-i\frac{p^{+}}{k^{+}}s}}{p^{+}(p^{+}+k^{+})}\,f(p^{+})
& =\underbrace{\int_{0}^{\infty}\!\!dp^{+}\,\frac{k^{+}\,e^{-i\frac{p^{+}}{k^{+}}s}}{p^{+}(p^{+}+k^{+})}\,\big[f(p^{+})-f(0)\big]}_{\text{plus term, phase kept in the subtraction}}\notag\\ & \quad
+\underbrace{f(0)\,\Big[\log\frac{k^{+}}{\Lambda}+\log\Big(1-\frac{k^{+}}{\vee}\Big)\Big]}_{\text{JIMWLK list, unchanged}}
+\underbrace{f(0)\,\Psi(s)}_{\text{extra finite term}}+{\cal O}\Big(\frac{\Lambda}{k^{+}},\frac{k^{+}}{\vee}\Big)\,,\notag
\end{align}
where $f(p^{+})$ is what multiplies the phase. At fixed coordinates the first and the third term add up to the author's plus-prescription term with $\Lambda\to0$; the third term is the part that grows like $-\log|k\cdot(z-y')|$ at large $z$. In Row III (first part) and Rows V, VI $f$ does not depend on $p^{+}$ and the plus term vanishes identically.
\subsection{Row I}
\begin{align}
     \frac{d^{3}N}{d^{3}k}  & =  -\,\frac{1}{(2\pi)^{3}}\,\frac{ig^{4}}{8\pi^{5}}\,\frac{1}{k^{+}} \int_{x,y,x',z,w,w^{\prime}} e^{-ik\cdot(w-w^{\prime})}\,\mathcal{K}^{i}(x'-w')\,\mathcal{K}^{j}(y-z)\,\log\left[\frac{(x-z)^{2}}{(x-w)^{2}}\right]\,\mathcal{T}^{ij}(x,z,w)\notag \\
& \times \Big[U^{ab'}(x')\,U^{c'd'}(y)\,U^{c'b}(z)\,\rho^{b'}(x')\rho^{d'}(y)-U^{d'b}(z)\,U^{d'e'}(y)U^{ac'}(w')\,\rho^{c'}(x')\rho^{e'}(y)\Big]\notag \\ & \times\Big[f^{cbd}\,U^{ac}(w)\,\rho^{d}(x)-f^{acd}\, U^{cb}(z)\,U^{de}(x)\,\rho^{e}(x)\Big]
\end{align}
\subsection{Row II}
\begin{align}
     \frac{d^{3}N}{d^{3}k}  & =  -\,\frac{1}{(2\pi)^{3}}\,\frac{ig^{4}}{8\pi^{5}}\,\frac{1}{k^{+}} \int_{x,y,x',z,w,w^{\prime}} e^{-ik\cdot(w-w^{\prime})}\,\mathcal{K}^{i}(x'-w')\,\mathcal{K}^{j}(y-z)\,\log\left[\frac{(x-z)^{2}}{(x-w)^{2}}\right]\,\mathcal{T}^{ij}(x,z,w)\notag \\
& \times \Big[U^{c'd'}(y)\,U^{ab'}(x')\,U^{c'b}(z)\,\rho^{d'}(y)\rho^{b'}(x')-U^{ab'}(x')\,\rho^{b}(y)\rho^{b'}(x')\Big]\Big[f^{cbd}\,U^{ac}(w)\,\rho^{d}(x)-f^{acd}\, U^{cb}(z)\,U^{de}(x)\,\rho^{e}(x)\Big]
\end{align}
{\color{gray}\small From Rows I, II of the checked row; same as the limits file. The $p^{+}$ integrand and the kernel are those of $N_{1a}$ of Group II Row I (prefactor $\times2$): same numerical check ($9.0\times10^{-4}\dots9.1\times10^{-10}$), same integrability, no $1/z^{2}$ tail.}
\subsection{Row III}
\noindent Colour factors of this row (the index $a$ is contracted with $f^{c'd'a}$ in the first part):
\begin{align}
\mathcal{L}^{c'd'b}(y') & \equiv U^{c'e'}(y')\,U^{d'b}(z)\,\rho^{e'}(x')-U^{d'b}(z)\,U^{c'e'}(x')\,\rho^{e'}(x')\,,\notag\\
\mathcal{R}_{1}^{ab} & \equiv U^{ac}(w)\, \{\rho^{c}(y),\rho^{b}(x)\}-U^{cb}(z)\,U^{ad}(y)\,U^{ce}(x)\,\{\rho^{d}(y),\rho^{e}(x)\}\,,\notag\\
\mathcal{R}_{2}^{ab} & \equiv f^{cbd}\,U^{ac}(w)\,\rho^{d}(x) -f^{acd}\,U^{cb}(z)\,U^{de}(x)\,\rho^{e}(x)\,;\notag
\end{align}
$\mathcal{L}^{c'd'b}(z)$ is $\mathcal{L}^{c'd'b}(y')$ with $U^{c'e'}(y')\to U^{c'e'}(z)$.
\subsubsection{First part}
\begin{align}
\frac{d^{3}N_{1}^{\widehat{\mathcal{E}}}}{d^{3}k} & = \frac{1}{(2\pi)^{3}}\, \frac{ig^{4}\,f^{c'd'a}}{16\pi^{5}}\, \frac{1}{k^{+}}\int_{w,x',y',z,y,x} e^{-ik\cdot (w-y')}\;\widehat{\mathcal{E}}\big(k\cdot(z-y')\big)\,\Big[\mathcal{K}(y'-z)\cdot\mathcal{K}(x'-y')\;\mathcal{K}(y-w)\cdot\mathcal{K}(x-z)\notag \\ & \qquad-\mathcal{K}(y'-z)\cdot\mathcal{K}(x-z)\;\mathcal{K}(x'-y')\cdot\mathcal{K}(y-w)\Big]\;\mathcal{L}^{c'd'b}(y')\;\mathcal{R}_{1}^{ab}\\
\frac{d^{3}N_{1}^{\Psi}}{d^{3}k} & = -\,\frac{1}{(2\pi)^{3}}\, \frac{ig^{4}\,f^{c'd'a}}{16\pi^{5}}\, \frac{1}{k^{+}}\int_{w,x',y',z,y,x} e^{-ik\cdot (w-y')}\;\Psi\big(k\cdot(z-y')\big)\;\mathcal{K}(y'-z)\cdot\mathcal{K}(x-z)\;\mathcal{K}(x'-y')\cdot\mathcal{K}(y-w)\;\mathcal{L}^{c'd'b}(y')\;\mathcal{R}_{1}^{ab}\\
\frac{d^{3}N_{1}^{\rm sub}}{d^{3}k} & = -\,\frac{1}{(2\pi)^{3}}\, \frac{ig^{4}\,f^{c'd'a}}{8\pi^{4}}\, \frac{1}{k^{+\,2}}\int_{0}^{\infty}\frac{dp^{+}}{2\pi}\int_{w,x',y',z,y,x} e^{-ik\cdot (w-y')}\, e^{-ik\cdot (z-y')\frac{p^{+}}{k^{+}}}\;\mathcal{K}(y'-z)\cdot\mathcal{K}(y-w)\;\mathcal{K}^{j}(x-z)\notag \\ & \qquad\times\Big[\mathcal{K}^{j}(x'-y')\,\mathcal{L}^{c'd'b}(y')-\mathcal{K}^{j}(x'-z)\,\mathcal{L}^{c'd'b}(z)\Big]\;\mathcal{R}_{1}^{ab}\qquad\text{($y'$ integral before the $p^{+}$ integral)}\notag\\
& = -\,\frac{1}{(2\pi)^{3}}\, \frac{ig^{4}\,f^{c'd'a}}{16\pi^{5}}\, \frac{1}{k^{+}}\int_{w,x',y',z,y,x} e^{-ik\cdot (w-y')}\,\Big[\pi\,\delta(s)-i\,{\rm P}\frac{1}{s}\Big]_{s=k\cdot(z-y')}\;\mathcal{K}(y'-z)\cdot\mathcal{K}(y-w)\;\mathcal{K}^{j}(x-z)\notag \\ & \qquad\times\Big[\mathcal{K}^{j}(x'-y')\,\mathcal{L}^{c'd'b}(y')-\mathcal{K}^{j}(x'-z)\,\mathcal{L}^{c'd'b}(z)\Big]\;\mathcal{R}_{1}^{ab}
\end{align}
In the distributional form the component of $y'-z$ perpendicular to $k$ is integrated first (symmetrically at large distance), then the component along $k$ with the principal value; the two terms of the last bracket stay together. The first and the second term are the terms with $\frac{1}{p^{+}+k^{+}}$ and the extra finite term $f(0)\Psi$; their part with the structure $\mathcal{K}(y'-z)\cdot\mathcal{K}(x-z)\,\mathcal{K}(x'-y')\cdot\mathcal{K}(y-w)$ adds up to the coefficient $+\big[\gamma_{E}+\log|s|+\frac{i\pi}{2}{\rm sgn}(s)\big]$.\par
{\color{gray}\small From the first part of Row III of the checked row (form with $\big[\frac{\delta_{k'm}\delta_{ij}}{p^{+}+k^{+}}-\frac{1}{p^{+}}\delta_{jm}\delta_{ik'}-\frac{1}{k^{+}}\delta_{im}\delta_{jk'}\big]$, range $\Lambda<p^{+}<\vee-k^{+}$) with the author's split (\texttt{Other\_terms\_checked}). Checked at fixed $s$ ($s=2.3$, $-0.37$): $\int_{\Lambda/k^{+}}^{\vee/k^{+}-1}\frac{dt}{1+t}e^{-its}-\widehat{\mathcal{E}}(s)$: $5.6\times10^{-3}\dots5.7\times10^{-6}$, and $\int\frac{dt}{t}e^{-its}-\log\frac{k^{+}}{\Lambda}-\log(1-\frac{k^{+}}{\vee})-[\Psi+\widehat{\mathcal{E}}](s)$: $2.8\times10^{-2}\dots2.7\times10^{-5}$ for $\frac{\Lambda}{k^{+}}=\frac{k^{+}}{\vee}=10^{-2}\dots10^{-5}$; closed form of $\Psi$ against quadrature: 10 digits. $N_{1}^{\rm sub}$, with a smooth test function $b$ for $\mathcal{K}^{j}(x'-y')\mathcal{L}(y')$: $\int_{0}^{\vee-k^{+}}dp^{+}\int_{y'}e^{i(1+\frac{p^{+}}{k^{+}})k\cdot(y'-z)}\mathcal{K}^{m}(y'-z)[b(y')-b(z)]$ differs from its value at $\vee=\infty$ by $0.64$, $7.0\times10^{-2}$, $7.0\times10^{-3}$, $7.0\times10^{-4}$ at $\frac{\vee}{k^{+}}=10,\dots,10^{4}$ (value $-3.78-2.48i$); the distributional form gives the same number to $7\times10^{-8}$; without the subtraction the integral grows like $2\pi i\frac{k^{m}}{k^{2}}\log\vee$ (the BFKL term).}\par
\noindent Integrability. $N_{1}^{\widehat{\mathcal{E}}}$: products of four kernels $1/r$, $\mathcal{K}(x'-y')[U(y')-U(x')]$ bounded; logarithm on the line $k\cdot(z-y')=0$; large $z$: $\widehat{\mathcal{E}}\simeq\frac{1}{ik\cdot z}$, hence $1/|z|^{3}$. $N_{1}^{\rm sub}$: the bracket vanishes at $y'\to z$; after the $y'$ integral the $p^{+}$ integrand is ${\cal O}(k^{+2}/p^{+2})$; large $z$: phase $e^{-ik\cdot(w-z)}$, no tail without phase. $N_{1}^{\Psi}$: integrable at all coincidences, {\color{magenta}but not at large $z$: tail $\frac{1}{z^{2}}\big[-\gamma_{E}-\log|k\cdot z|-\frac{i\pi}{2}{\rm sgn}(k\cdot z)\big]$ times the coefficient of the table in Sec.~\ref{sec:tails}; see OPEN II/III-2, 3.}
\subsubsection{Second part}
\noindent Notation: $a\equiv x-w$, $b\equiv x-z$, $c\equiv z-w$, $A\equiv a^{2}$, $B\equiv b^{2}$, $C\equiv c^{2}$, $D_{p}\equiv p^{+}(x-z)^{2}+k^{+}(x-w)^{2}$, $\mathcal{L}_{2}^{ab}(y')\equiv f^{c'd'a}\mathcal{L}^{c'd'b}(y')$,
\begin{align}
\widehat{\mathcal{J}} & \equiv\frac{a\cdot\mathcal{K}(x'-y')\;\,c\cdot\mathcal{K}(y'-z)}{C}+\frac{b\cdot\mathcal{K}(y'-z)\;\,a\cdot\mathcal{K}(x'-y')}{2\,B}\,,\notag\\
\mathcal{V} & \equiv\frac{b\cdot\mathcal{K}(y'-z)\;\,c\cdot\mathcal{K}(x'-y')}{C}-\frac{b\cdot\mathcal{K}(y'-z)\;\,a\cdot\mathcal{K}(x'-y')}{2\,A}+\frac{a\cdot\mathcal{K}(y'-z)\;\,c\cdot\mathcal{K}(x'-y')}{C}+\frac{b\cdot\mathcal{K}(x'-y')\;\,a\cdot\mathcal{K}(y'-z)}{2\,B}\,,\notag\\
\mathcal{W}^{k'm} & \equiv\frac{b^{k'}\,c^{m}}{C}-\frac{b^{k'}\,a^{m}}{2\,A}\,.\notag
\end{align}
\begin{align}
\frac{d^{3}N_{2}^{\rm reg}}{d^{3}k} & = \frac{1}{(2\pi)^{3}}\,\frac{g^{4}}{4\pi^{4}}\int_{0}^{\infty}\frac{dp^{+}}{2\pi}\int_{w,x',y',z,x} e^{-ik\cdot (w-y')}\,e^{-ik\cdot(z-y')\frac{p^{+}}{k^{+}}}\;\frac{1}{D_{p}}\,\Bigg\{\mathcal{K}(y'-z)\cdot\mathcal{K}(x'-y')\,\bigg[\frac{d_{\perp}\,p^{+}\,(B-A)}{2\,C\,(p^{+}+k^{+})^{2}}\notag \\ & \quad+\frac{p^{+}}{k^{+}\,(p^{+}+k^{+})}\Big(\frac{b\cdot c}{C} -\frac{a\cdot b}{2\,A}-\frac{B-A}{2\,C}\Big) +\frac{1}{p^{+}+k^{+}}\Big(\frac{a\cdot c}{C}+\frac{a\cdot b}{2\,B} -\frac{B-A}{2\,C}\Big)\bigg]\notag \\ & \quad-\bigg[\frac{k^{+}}{p^{+}(p^{+}+k^{+})}\Big(1-\frac{D_{p}}{k^{+}A}\Big)+\frac{1}{p^{+}+k^{+}}\bigg]\,\widehat{\mathcal{J}}\;-\;\frac{1}{k^{+}}\,\mathcal{V}\Bigg\}\;\mathcal{L}_{2}^{ab}(y')\;\mathcal{R}_{2}^{ab}\,,\qquad d_{\perp}=2\\
\frac{d^{3}N_{2}^{\Psi}}{d^{3}k} & = -\,\frac{1}{(2\pi)^{3}}\,\frac{g^{4}}{8\pi^{5}}\,\frac{1}{k^{+}}\int_{w,x',y',z,x} e^{-ik\cdot (w-y')}\;\Psi\big(k\cdot(z-y')\big)\;\frac{\widehat{\mathcal{J}}}{(x-w)^{2}}\;\mathcal{L}_{2}^{ab}(y')\;\mathcal{R}_{2}^{ab}\\
\frac{d^{3}N_{2}^{\rm sub}}{d^{3}k} & = -\,\frac{1}{(2\pi)^{3}}\,\frac{g^{4}}{4\pi^{4}}\,\frac{1}{k^{+\,2}}\int_{0}^{\infty}\frac{dp^{+}}{2\pi}\int_{w,x',y',z,x} e^{-ik\cdot (w-y')}\,e^{-ik\cdot(z-y')\frac{p^{+}}{k^{+}}}\;\frac{p^{+}}{D_{p}}\;\mathcal{K}^{m}(y'-z)\,\mathcal{W}^{k'm}\notag \\ & \qquad\times\Big[\mathcal{K}^{k'}(x'-y')\,\mathcal{L}_{2}^{ab}(y')-\mathcal{K}^{k'}(x'-z)\,\mathcal{L}_{2}^{ab}(z)\Big]\;\mathcal{R}_{2}^{ab}\qquad\text{($y'$ integral before the $p^{+}$ integral)}\\
\frac{d^{3}N_{2}^{\log}}{d^{3}k} & = \frac{1}{(2\pi)^{3}}\,\frac{g^{4}}{8\pi^{5}}\,\frac{1}{k^{+}}\int_{w,x',w',z,x} e^{-ik\cdot (w-w')}\;\frac{(x-w)^{2}}{(x-z)^{2}\,\left[(x-w)^{2}-(x-z)^{2}\right]}\,\log\left[\frac{(x-w)^{2}}{(x-z)^{2}}\right]\notag \\ & \qquad\times\mathcal{K}^{m}(w'-z)\,\mathcal{K}^{k'}(x'-z)\,\mathcal{W}^{k'm}\;\mathcal{L}_{2}^{ab}(z)\;\mathcal{R}_{2}^{ab}
\end{align}
In $N_{2}^{\rm reg}$ the bracket in front of $\widehat{\mathcal{J}}$ is the plus term with the phase kept in the subtraction plus the $\frac{1}{p^{+}+k^{+}}$ term; it equals $-\frac{B-A}{(p^{+}+k^{+})\,A}$, so that no $\frac{1}{p^{+}}$ is left. The whole brace of $N_{2}^{\rm reg}$ has to stay under one $p^{+}$ and one $z$ integral: at large $z$ and small $p^{+}$ the $\widehat{\mathcal{J}}$ term and the term $+\frac{1}{k^{+}}\frac{b\cdot\mathcal{K}(y'-z)\,a\cdot\mathcal{K}(x'-y')}{2A}$ of $-\frac{1}{k^{+}}\mathcal{V}$ are each $\propto\frac{1}{p^{+}z^{2}+k^{+}A}$ (their $z$ integrals grow like $\frac{1}{p^{+}}\log\frac{1}{p^{+}}$) and add up to $\frac{a\cdot\mathcal{K}(x'-y')}{2A}\,\frac{p^{+}}{k^{+}(p^{+}+k^{+})}\,\frac{1}{p^{+}z^{2}+k^{+}A}$, which is integrable. $N_{2}^{\log}$ is what is left of the term added back for the subtraction in $N_{2}^{\rm sub}$ after its $\frac{1}{(x-z)^{2}}\log\frac{\vee}{k^{+}}$ has gone to the BFKL list.\par
{\color{gray}\small From the second part of Row III of the checked row (form after the $w'$ integration) and the limits file. Checked at a random point ($s=0.081$), all terms except the constant-kernel one: [$p^{+}$ integral from $\Lambda$ to $\vee-k^{+}$] $-\big[\log\frac{k^{+}}{\Lambda}+\log(1-\frac{k^{+}}{\vee})\big]f(0)$ minus [integrand of $N_{2}^{\rm reg}$ $+$ integrand of $N_{2}^{\Psi}$] $=0.14$, $1.9\times10^{-2}$, $1.6\times10^{-3}$, $1.5\times10^{-4}$ for $\frac{\Lambda}{k^{+}}=\frac{k^{+}}{\vee}=10^{-2}\dots10^{-5}$ (value $-3.311+1.837i$). $N_{2}^{\rm sub}+N_{2}^{\log}$ with a test function: $\int_{0}^{\vee-k^{+}}dp^{+}\frac{p^{+}}{D_{p}}\int_{y'}\dots$ minus the added-back logarithms differs from the limit by $0.33\dots3.7\times10^{-4}$ at $\frac{\vee}{k^{+}}=10\dots10^{4}$. Large $z$, small $p^{+}$: the combination quoted above is reproduced to $2\times10^{-4}$ at $|z|=10^{5}$, $p^{+}=0.1k^{+}$.}\par
\noindent Integrability. $N_{2}^{\rm reg}$ ($p^{+}$ integrated at fixed points, numerically): $x\to w$, $x\to z$, $z\to w$, $x'\to y'$: $1/r$; $y'\to z$: $\log r/r$; large $z$: $R^{3}|N_{2}^{\rm reg}|=3.5$, $5.0$, $6.7$, $8.4$ at $|z|=R=10^{2},\dots,10^{5}$, i.e. $\log R/R^{3}$: no $1/z^{2}$ tail; $d_{\perp}=2$ (no pole multiplies it). $N_{2}^{\rm sub}$, $N_{2}^{\log}$: as in the first part and as in the limits file ($x\to z$: $\log/|x-z|$; regular on $(x-z)^{2}=(x-w)^{2}$). $N_{2}^{\Psi}$: {\color{magenta}tail $\frac{1}{z^{2}}\big[-\gamma_{E}-\log|k\cdot z|-\frac{i\pi}{2}{\rm sgn}(k\cdot z)\big]$ times the coefficient of the table; OPEN II/III-2, 3.}
\subsection{Row IV}
\textsc{Only 4 $\rho$'s}
\subsection{Row V}
\noindent $\mathcal{L}_{\rm V}^{c'}(y')\equiv U^{c'e'}(y')\,\rho^{e'}(x')-U^{c'e'}(x')\,\rho^{e'}(x')$, $\;\mathcal{R}_{\rm V}^{ac}\equiv U^{ce}(x)\,U^{ad}(y)\,\rho^{e}(x)\rho^{d}(y)-U^{ad}(w)\,U^{ce}(x)\,\rho^{e}(x)\rho^{d}(y)$.
\begin{align}
\frac{d^{3}N^{\widehat{\mathcal{E}}}}{d^{3}k} & = \frac{1}{(2\pi)^{3}}\, \frac{ig^{4}\,f^{c'ca}}{8\pi^{5}}\, \frac{1}{k^{+}}\int_{w,y',z,y,x,x'} e^{-ik\cdot (w-y')}\;\widehat{\mathcal{E}}\big(k\cdot(z-y')\big)\,\Big[\mathcal{K}(y'-z)\cdot\mathcal{K}(x'-y')\;\mathcal{K}(y-w)\cdot\mathcal{K}(x-z)\notag \\ & \qquad-\mathcal{K}(y'-z)\cdot\mathcal{K}(x-z)\;\mathcal{K}(x'-y')\cdot\mathcal{K}(y-w)\Big]\;\mathcal{L}_{\rm V}^{c'}(y')\;\mathcal{R}_{\rm V}^{ac}\\
\frac{d^{3}N^{\Psi}}{d^{3}k} & = -\,\frac{1}{(2\pi)^{3}}\, \frac{ig^{4}\,f^{c'ca}}{8\pi^{5}}\, \frac{1}{k^{+}}\int_{w,y',z,y,x,x'} e^{-ik\cdot (w-y')}\;\Psi\big(k\cdot(z-y')\big)\;\mathcal{K}(y'-z)\cdot\mathcal{K}(x-z)\;\mathcal{K}(x'-y')\cdot\mathcal{K}(y-w)\;\mathcal{L}_{\rm V}^{c'}(y')\;\mathcal{R}_{\rm V}^{ac}\\
\frac{d^{3}N^{\rm sub}}{d^{3}k} & = -\,\frac{1}{(2\pi)^{3}}\, \frac{ig^{4}\,f^{c'ca}}{8\pi^{5}}\, \frac{1}{k^{+}}\int_{w,y',z,y,x,x'} e^{-ik\cdot (w-y')}\,\Big[\pi\,\delta(s)-i\,{\rm P}\frac{1}{s}\Big]_{s=k\cdot(z-y')}\;\mathcal{K}(y'-z)\cdot\mathcal{K}(y-w)\;\mathcal{K}^{j}(x-z)\notag \\ & \qquad\times\Big[\mathcal{K}^{j}(x'-y')\,\mathcal{L}_{\rm V}^{c'}(y')-\mathcal{K}^{j}(x'-z)\,\mathcal{L}_{\rm V}^{c'}(z)\Big]\;\mathcal{R}_{\rm V}^{ac}
\end{align}
\subsection{Row VI}
\noindent $\mathcal{R}_{\rm VI}^{ab}\equiv U^{cb}(z)\,U^{ad}(y)\,U^{ce}(x)\,\rho^{d}(y)\,\rho^{e}(x)-U^{ad}(y)\,\rho^{d}(y)\,\rho^{b}(x)$; $\mathcal{L}^{c'd'b}$ as in Row III.
\begin{align}
\frac{d^{3}N^{\widehat{\mathcal{E}}}}{d^{3}k} & = \frac{1}{(2\pi)^{3}}\, \frac{ig^{4}\,f^{c'd'a}}{8\pi^{5}}\, \frac{1}{k^{+}}\int_{w,y',z,y,x,x'} e^{-ik\cdot (w-y')}\;\widehat{\mathcal{E}}\big(k\cdot(z-y')\big)\,\Big[\mathcal{K}(y'-z)\cdot\mathcal{K}(x'-y')\;\mathcal{K}(y-w)\cdot\mathcal{K}(x-z)\notag \\ & \qquad-\mathcal{K}(y'-z)\cdot\mathcal{K}(x-z)\;\mathcal{K}(x'-y')\cdot\mathcal{K}(y-w)\Big]\;\mathcal{L}^{c'd'b}(y')\;\mathcal{R}_{\rm VI}^{ab}\\
\frac{d^{3}N^{\Psi}}{d^{3}k} & = -\,\frac{1}{(2\pi)^{3}}\, \frac{ig^{4}\,f^{c'd'a}}{8\pi^{5}}\, \frac{1}{k^{+}}\int_{w,y',z,y,x,x'} e^{-ik\cdot (w-y')}\;\Psi\big(k\cdot(z-y')\big)\;\mathcal{K}(y'-z)\cdot\mathcal{K}(x-z)\;\mathcal{K}(x'-y')\cdot\mathcal{K}(y-w)\;\mathcal{L}^{c'd'b}(y')\;\mathcal{R}_{\rm VI}^{ab}\\
\frac{d^{3}N^{\rm sub}}{d^{3}k} & = -\,\frac{1}{(2\pi)^{3}}\, \frac{ig^{4}\,f^{c'd'a}}{8\pi^{5}}\, \frac{1}{k^{+}}\int_{w,y',z,y,x,x'} e^{-ik\cdot (w-y')}\,\Big[\pi\,\delta(s)-i\,{\rm P}\frac{1}{s}\Big]_{s=k\cdot(z-y')}\;\mathcal{K}(y'-z)\cdot\mathcal{K}(y-w)\;\mathcal{K}^{j}(x-z)\notag \\ & \qquad\times\Big[\mathcal{K}^{j}(x'-y')\,\mathcal{L}^{c'd'b}(y')-\mathcal{K}^{j}(x'-z)\,\mathcal{L}^{c'd'b}(z)\Big]\;\mathcal{R}_{\rm VI}^{ab}
\end{align}
{\color{gray}\small Rows V, VI: from the checked rows (same kernels as Row III, first part; prefactor $\frac{ig^{4}f}{4\pi^{4}k^{+}}$ instead of $\frac{ig^{4}f}{8\pi^{4}k^{+}}$); the numerical checks of Row III, first part, apply (they concern the functions of $s$ only). $N^{\rm sub}$ is written in the distributional form; equivalently $\big[\pi\delta(s)-i{\rm P}\frac1s\big]\to\frac{1}{k^{+}}\int_{0}^{\infty}dp^{+}e^{-i\frac{p^{+}}{k^{+}}s}$ with the $y'$ integral done first. Integrability as in Row III, first part. In Row V the colour factor does not contain $U(z)$.} {\color{magenta}$N^{\Psi}$ of Rows V and VI: tail as in Row III; OPEN II/III-2, 3.}
\subsection{Large-distance tails of the finite terms of Groups II and III}\label{sec:tails}
\noindent Finite terms without a $1/z^{2}$ tail (at fixed $w$, $w'$): Group II Row I ($N_{1a}$, $N_{1b}$; the sum of $\mathfrak{K}_{2a},\dots,\mathfrak{K}_{2f}$), Group III Rows I, II, all $N^{\widehat{\mathcal{E}}}$, $N^{\rm sub}$, $N_{2}^{\rm reg}$, $N_{2}^{\log}$. Finite terms with a tail: Group II Row IV (OPEN II/III-1; variable $Y$, no recoil), and the extra terms $N^{\Psi}$ of Group III. For the latter, with $y'\to w'$,
\begin{align}
\text{integrand of }N^{\Psi}\;\xrightarrow{|z|\to\infty}\;c_{0}\;e^{-ik\cdot(w-w')}\;\frac{1}{z^{2}}\,\Big[-\gamma_{E}-\log|k\cdot z|-\frac{i\pi}{2}\,{\rm sgn}(k\cdot z)\Big]\;\mathsf{T}\,,\qquad c_{0}\equiv\frac{1}{(2\pi)^{3}}\frac{g^{4}}{16\pi^{5}}\frac{1}{k^{+}}\,,\notag
\end{align}
with the coefficients $\mathsf{T}$ of the table (they are the tails of the JIMWLK-list coefficients $f(0)$ of these rows). The terms written above have the recoil of the measured gluon on the conjugate side ($w'$; phase $e^{-ik\cdot(z-w')p^{+}/k^{+}}$); their complex conjugates have it on the side of $w$, with $\mathsf{T}\to\mathsf{T}^{\dagger}$ and the square bracket replaced by its complex conjugate. $\Delta^{c'e'}\equiv U^{c'e'}(w')-U^{c'e'}(x')$; $U(z)$ is the Wilson line at the far point; in the $z$-independent entries of Rows III and VI $U^{d'b}(z)U^{cb}(z)=\delta^{d'c}$ has been used (these are the terms of second degree in $U(z)$).
\begin{center}
\renewcommand{\arraystretch}{1.45}
\small\begin{tabular}{|l|l|l|l|}
\hline
term (side $w'$) & kernel & degree in $U(z)$ & coefficient $\mathsf{T}$ (colour factor) \\ \hline
III, Row III, $N_{1}^{\Psi}$ & $\mathcal{K}(x'-w')\!\cdot\!\mathcal{K}(y-w)$ & linear & $-i f^{c'd'a}\,\Delta^{c'e'}U^{d'b}(z)\,U^{ac}(w)\,\rho^{e'}(x')\,\{\rho^{c}(y),\rho^{b}(x)\}$ \\
 & & $z$-independent & $+i f^{c'ca}\,\Delta^{c'e'}U^{ad}(y)\,U^{ce}(x)\,\rho^{e'}(x')\,\{\rho^{d}(y),\rho^{e}(x)\}$ \\ \hline
III, Row III, $N_{2}^{\Psi}$ & $\mathcal{K}(x'-w')\!\cdot\!\mathcal{K}(x-w)$ & linear & $+f^{c'd'a}f^{cbd}\,\Delta^{c'e'}U^{d'b}(z)\,U^{ac}(w)\,\rho^{e'}(x')\,\rho^{d}(x)$ \\
 & & $z$-independent & $+N_{c}\,\Delta^{de'}U^{de}(x)\,\rho^{e'}(x')\,\rho^{e}(x)$ \\ \hline
III, Row V, $N^{\Psi}$ & $\mathcal{K}(x'-w')\!\cdot\!\mathcal{K}(y-w)$ & linear & $0$ \\
 & & $z$-independent & $-2i f^{c'ca}\,\Delta^{c'e'}\rho^{e'}(x')\,\big[U^{ce}(x)\,U^{ad}(y)-U^{ad}(w)\,U^{ce}(x)\big]\,\rho^{e}(x)\,\rho^{d}(y)$ \\ \hline
III, Row VI, $N^{\Psi}$ & $\mathcal{K}(x'-w')\!\cdot\!\mathcal{K}(y-w)$ & linear & $+2i f^{c'd'a}\,\Delta^{c'e'}U^{d'b}(z)\,U^{ad}(y)\,\rho^{e'}(x')\,\rho^{d}(y)\,\rho^{b}(x)$ \\
 & & $z$-independent & $-2i f^{c'ca}\,\Delta^{c'e'}U^{ad}(y)\,U^{ce}(x)\,\rho^{e'}(x')\,\rho^{d}(y)\,\rho^{e}(x)$ \\ \hline
II, Row IV, (i)+(ii)+(iii) & (variable $Y$) & --- & $2N_{c}\,\mathcal{C}\,\frac{4\pi}{k^{2}}\,\frac{1}{Y^{2}}\big[-\frac{11}{6}+\frac16\cos2\varphi\big]\,G_{\infty}(Y)$ \\ \hline
\end{tabular}
\end{center}
(The last line is the complete tail of OPEN II/III-1: no recoil, no factor $c_{0}$, no logarithm. For $N_{1}^{\Psi}$ and Rows V, VI the kernel does not depend on $x$; in the entries linear in $U(z)$ of $N_{1}^{\Psi}$ and of Row VI the point $x$ enters only through $\int_{x}\rho^{b}(x)$.) Maximum norms of the symmetrised two-$\rho$ parts at two sites, $SU(3)$, per side [$w'$; $w$], by degree in $U(z)$ (linear; $z$-independent): Row III first part $[1.66,\,0.83;\ 1.33,\,0.39]$, Row III second part $[0.77,\,1.62;\ 0.55,\,1.80]$, Row V $[0,\,1.60;\ 0,\,1.60]$, Row VI $[2.23,\,2.45;\ 2.27,\,1.88]$.\par
{\color{gray}\small The entries are the large-$z$ limits of the kernels of $N^{\Psi}$ ($\mathcal{K}(y'-z)\cdot\mathcal{K}(x-z)\to\frac{1}{z^{2}}$, $\widehat{\mathcal{J}}/A\to-\frac{\mathcal{K}(x'-y')\cdot\mathcal{K}(x-w)}{2z^{2}}$). Checked with the lattice engine of the earlier tail analysis ($SU(2)$, $SU(3)$, Weyl reduction, $|z|=2\times10^{4}$): the entries agree with the numerical tails of the rows to $2\times10^{-8}$.}\par\smallskip
\noindent\textbf{Sums.} (a) The $z$-independent entries of Row III (both parts) and Row VI cancel identically on each side: with $[\rho^{e}(x),\rho^{d}(y)]=if^{edh}\rho^{h}(x)\delta^{(2)}(x-y)$, [Row III first part] $+$ [Row VI] $=-N_{c}\,\mathcal{K}(x'-w')\cdot\mathcal{K}(x-w)\,\Delta^{de'}U^{de}(x)\rho^{e'}(x')\rho^{e}(x)$, minus the entry of the second part [numerically: two-$\rho$ part $2\times10^{-8}$, three-$\rho$ part $10^{-15}$, single entries of order $1$--$2$]. (b) What is left of Group III on each side is the $z$-independent entry of Row V and the three entries linear in $U(z)$. Tail of the Group III part [$SU(3)$, norms as above]: side $w'$: $z$-independent $1.603$ (Row V), linear $3.188$; side $w$: $1.603$ and $3.188$; the linear parts of the two sides together: $4.844$. (c) Adding the recoil rows of Group I, with the weights of \texttt{Cancellation\_of\_flagged\_terms} (Row V first part, $N_{2a}$, $N_{2b\,\rm rest}$: weight $1$; Row III, term with $\frac12K(x,y,z)$: weight $2$; Row III, $z\to w$ term in the soft-position form, $U(z)\to U(w)$: weight $1$): the $z$-independent tails cancel on each side separately [$1\times10^{-8}$], the tails linear in $U(z)$ cancel only between the two sides [each side $2.366$; sum $1\times10^{-8}$] ($SU(2)$: the same pattern).\par\smallskip
\noindent{\color{magenta}\fbox{OPEN II/III-2} \textbf{The extra terms $N^{\Psi}$ are not integrable over $z$ one by one.} Each of $N_{1}^{\Psi}$, $N_{2}^{\Psi}$ (Row III), $N^{\Psi}$ (Rows V, VI) diverges at large $z$ like $-\pi\,\mathsf{T}\log^{2}$ of the large-distance cut-off [model kernel $\frac{1}{z^{2}+1}$ with $\Lambda$ kept: second difference in $\log\frac{k^{+}}{\Lambda}$ equal to $-2\pi$ to $10^{-4}$]. They are defined only in the sum, under one $z$ integral, of: the four $N^{\Psi}$ terms of Group III with their complex conjugates, and the extra terms generated in the same way (phase kept in the subtraction) by Group I Row V (first part, $N_{2a}$, $N_{2b\,\rm rest}$) and by the two flagged terms of Group I Row III (weights above). These Group I terms are not part of this section; their tail function has to be $-\gamma_{E}-\log|k\cdot z|\mp\frac{i\pi}{2}{\rm sgn}(k\cdot z)$ as well (this is the large-$|s|$ form of $\int_{\Lambda}\frac{dq^{+}}{q^{+}}e^{\mp i\frac{q^{+}}{k^{+}}s}-\log\frac{k^{+}}{\Lambda}$ for any upper limit of order $k^{+}$, but it was not checked on the Group I rows here). For a combination without $1/z^{2}$ tail the limit $\Lambda\to0$ after the $z$ integration exists and is the integral of the pointwise limit [model: $\int_{z}\big[\frac{1}{z^{2}+1}-\frac{1}{(z-c)^{2}+\rho^{2}}\big]\Psi_{\Lambda}(k\cdot z)$ differs from its limit by $0.25$, $3.4\times10^{-2}$, $4.3\times10^{-3}$, $5.2\times10^{-4}$ at $\frac{\Lambda}{k^{+}}=10^{-2}\dots10^{-5}$, value $3.09+6.40i$]; this was tested with model kernels, not with the full rows.}\par\smallskip
\noindent{\color{magenta}\fbox{OPEN II/III-3} \textbf{Tails linear in $U(z)$.} They cancel only between the terms with the recoil on $w'$ (tail function $-\gamma_{E}-\log|k\cdot z|-\frac{i\pi}{2}{\rm sgn}(k\cdot z)$) and those with the recoil on $w$ (complex-conjugate function). The real parts cancel; left is $-i\pi\,{\rm sgn}(k\cdot z)\,\frac{1}{z^{2}}\,\mathsf{T}^{(1)}_{w'}[U(z)]$ ($\mathsf{T}^{(1)}_{w'}$: sum of the linear entries of all recoil rows on the side $w'$, Group I included), which integrates to zero at large $z$ only if $U(z)$ becomes independent of the direction of $z$ (or even under $z\to-z$) at large distance. Nothing is assumed about $U$ at large distance, so this is not decided.}\par\smallskip
\noindent{\color{magenta}\fbox{OPEN II/III-4} \textbf{Not among the finite terms.} (a) The $\log\big(1-\frac{k^{+}}{\vee}\big)$ terms of all rows: they are $\log(1-\frac{k^{+}}{\vee})$ times the JIMWLK structures of the rows, whose $z$ integrals diverge logarithmically for a general projectile; they vanish for $\vee\to\infty$ only with the large-distance regulator in place and are kept with the $\log\frac{k^{+}}{\Lambda}$ terms. (b) The term $\log\frac{k^{+}}{\Lambda}\int\frac{d^{2}z}{z^{2}}\,C_{A}$ of the JIMWLK formula for a non-singlet pair of charges is untouched by anything in this section. (c) The poles $\frac{2}{\epsilon}$ of Row IV are not written.}\par\smallskip
\noindent\textbf{Not checked.} The full rows of Group III were not integrated numerically over $z$ with Wilson lines (the limits were checked pointwise, and after integration with test functions and model kernels); the tails were obtained at fixed $w$, $w'$ (regions where $w$, $w'$ go to infinity together with $z$: power counting only); in the lattice evaluation the three-$\rho$ rows are evaluated with two charge sites; the $\log^{2}$ plateaux of the single rows rest on the analytic large-$z$ form of the kernels.

\end{document}
